\documentclass[12pt,a4paper]{book}
% \usepackage{times}
\usepackage[a4paper, margin=1.25in]{geometry}
\usepackage[utf8]{inputenc}
\usepackage[english]{babel}
\selectlanguage{english}
\usepackage{amsmath,amsthm,amssymb,fge,mathrsfs}
\usepackage{mathtools}
% \usepackage{hyperref}
\usepackage[hidelinks, colorlinks = true, linkcolor = gray, citecolor = gray]{hyperref}
\usepackage{csquotes}

\usepackage{makeidx} %per fare l'indice
\usepackage{faktor} %per i quozienti

%tikz (pacchetto per i disegni)
\usepackage{tikz}
\usetikzlibrary{arrows}
\usetikzlibrary{tikzmark}
\usetikzlibrary{calc} %per poter fare i calcoli
\usetikzlibrary{arrows.meta} %per usare veri tipi di frecce
\usetikzlibrary{calc,patterns,angles,quotes} %per disegnare gli angoli
\usetikzlibrary{decorations} %per i grafici orientati (con sopra le frecce)
\usetikzlibrary{decorations.markings}
\usetikzlibrary{backgrounds} %per il colore sullo sfondo delle immagini
\usetikzlibrary{shapes.geometric} %per i triangoli negli alberi

\usepackage{pgfplots}
\pgfplotsset{compat=newest}
\usepgfplotslibrary{fillbetween} %per colorare le aree comprese tra due grafici

\usepackage{caption} %per la caption sotto alle immagini
\usepackage[tracking=true]{microtype} %per avvicinare il testo
\usepackage{enumitem} %per modificare agilmente gli elenchi
\usepackage{esvect} %per le freccette di vettore sui caratteri
\usepackage{stmaryrd} %fulmine per il simbolo dell'assurdo
\usepackage{relsize} %per fare i simboli più grandi
\usepackage{mathabx} %per il simbolo di dotminus
\usepackage{esint} %per l'integrale tagliato (fint)

\usepackage[customcolors]{hf-tikz} %per evidenziare pezzi di matrice

\usepackage{verbatim} %per commentare un blocco

\usepackage{multirow} %per la tabella
\usepackage{array} %per lo spazio nella tabella

%Libreria per gli alberi
\usepackage[linguistics]{forest}

%Librerie aggiunte
\usepackage{multicol} %per gli elenchi puntati su più colonne/righe
\usepackage{multirow}

%pacchetto per il prodotto interno
\usepackage{physics}

%Libreria per i subfiles

\usepackage{subfiles} %meglio caricarla per ultima

\usepackage{biblatex}
% \addbibresource{Bibliography.bib}
\addbibresource{Bibliography.bib}


% \addbibresource{biblio.bib}
%usiamo paragraph come nuova sottosottosottosezione 
\makeatletter
\renewcommand\paragraph{\@startsection{paragraph}{4}{\z@}%
            {-2.5ex\@plus-1ex \@minus -.25ex}%
            {1.25ex \@plus .25ex}%
            {\normalfont\normalsize\bfseries}}
\makeatother
\setcounter{secnumdepth}{4} % how many sectioning levels to assign numbers to
\setcounter{tocdepth}{4}    % how many sectioning levels to show in ToC

\newtheorem{theorem}{Theorem}
\numberwithin{theorem}{section}
\newtheorem*{theorem*}{Theorem}
\newtheorem{corollary}[theorem]{Corollary}
\newtheorem{proposition}[theorem]{Proposition}
\newtheorem{lemma}[theorem]{Lemma}
\newtheorem*{lemma*}{Lemma}

\newtheorem{definition}{Definition}
\numberwithin{definition}{section}

\newtheorem*{remark}{Remark}

\newtheorem{example}{Example}
\numberwithin{example}{section}

\newtheorem{exercise}{Exercise}
\numberwithin{exercise}{section}

\renewcommand{\proofname}{\proof}

\renewcommand{\thefootnote}{[\arabic{footnote}]}% Modify footnote globally

%%%simboli speciali
\newcommand{\N}{\mathbb{N}}
\newcommand{\Z}{\mathbb{Z}}
\newcommand{\Q}{\mathbb{Q}}
\newcommand{\R}{\mathbb{R}}
\newcommand{\C}{\mathbb{C}}

\newcommand{\CO}{\mathcal{C}}

\DeclareMathOperator{\Dom}{Dom}
\DeclareMathOperator{\Ima}{Im}

\DeclareMathOperator{\mcd}{MCD}
\DeclareMathOperator{\mcm}{mcm}
\DeclareMathOperator{\supp}{supp}
\DeclareMathOperator{\spanlin}{span}

%\DeclarePairedDelimiter\abs{\lvert}{\rvert}
%\DeclarePairedDelimiter\norm{\lVert}{\rVert}

%Necessario per far si che la dimensione del valore assoluto e della norma si adatti all'aromento passato
%\makeatletter
%\let\oldabs\abs
%\def\abs{\@ifstar{\oldabs}{\oldabs*}}
%
%\let\oldnorm\norm
%\def\norm{\@ifstar{\oldnorm}{\oldnorm*}}
%\makeatother

%rinomino il comando per il prodotto interno
\let\pint\braket
%comando per il prodotto scalare
\newcommand{\pscal}[2]{\left\langle #1, #2 \right\rangle}

%Spaziatura per i quantificatori
\let\oldforall\forall
\renewcommand{\forall}{\; \oldforall \;}
\let\oldexists\exists
\renewcommand{\exists}{\; \oldexists \;}

%freccia per la convergenza debole
\newcommand{\longrightharpoonup}{\xrightharpoonup{\phantom{AB}}}

\newcommand{\warrow}{\xrightharpoonup{\:\, w \:\, }}
\newcommand{\sarrow}{\overset{s}{\longrightarrow}}
\newcommand{\wastarrow}{\xrightharpoonup{ w^* }}


\setlength{\parindent}{0in} 
\title{Title}
\author{Andrea Scalenghe}
\date{March 2024}



\begin{document}
\thispagestyle{empty}

\centerline {\Large{\textsc{ UNIVERSIT\`A DEGLI STUDI DI TORINO}}}
\vskip 20 pt

\centerline {\Large{\textsc DIPARTIMENTO DI MATEMATICA GIUSEPPE PEANO}}

\vskip 20 pt

\centerline {{\textsc SCUOLA DI SCIENZE DELLA NATURA}}

\vskip 20 pt

\centerline {\Large{\textsc Corso di Laurea Magistrale in Matematica}}
\vskip 60 pt





%\begin{tabular}{ccc}
\centerline {\includegraphics[width=4cm]{logo_new_2022.png}}
%   \end{tabular}

\vskip 1.2cm

\centerline {\normalsize {Tesi di Laurea  Magistrale}} 

\vskip 0.7cm

\centerline {\Large {\bf TITOLO}}

\vskip 1.7cm

\noindent Relatrice: Susanna Terracini
\hfill  {Candidato: Andrea Scalenghe}

\noindent Co-relatrice: Elena Issoglio


\vskip 2.7cm


\centerline{ANNO ACCADEMICO 2023-2024}


\newpage

\tableofcontents

\newpage

\chapter{Introduction}

In this dissertation I will adrress optimal control theory both from a deterministic and stochastic point of view. 

\newpage

\chapter{Deterministic Optimal Control}

We now study deterministic optimal control. The system we aim to control is governed by differential equations. 

Short description of what is done in this chapter.

\section{Finite horizon}


Let us consider a finite interval $I=[t,t_1]\subset\R$ as the operating time of the system. At each time $s\in I$ the system is described by $x(s)\in O\subseteq \R^n$ and controlled by $u(s)\in U\subseteq\R^n$
called control space. The system is described by:

\begin{equation}\label{1-1-syst1}
    \begin{cases}
        \dot{x}(s) = f(s,x(s),u(s)) & s\in I \\
        x(t) = x
    \end{cases}
\end{equation}

For a given $x\in O$ and suitable $f:\overline{Q}\times U \rightarrow \R^m$, where $Q_0=[t,t_1)\times O$. That is we impose $f\in C(\overline{Q}\times U)$ and the existence of $K_{\rho}>0$ for all $\rho>0$:

\begin{equation}\label{1-1-lipsch}
    \abs{f(t,x,v)-f(t,y,v)} \leq K_{\rho}\abs{x-y}
\end{equation}

For all $t\in I$, $x,y\in O$ and $v\in U$ such that $\abs{v}\leq\rho$. Under this conditions the system \ref*{1-1-syst1} has a unique solution. 
Controls $u(\cdot)$ are assumed to be in the set $L^{\infty}\left([t,t_1];U\right)$. We will soon specify more about the set of controls.

We have described a control problem. The concept of optimality is related some value function, specified by payoffs (or costs) associated to the system's states.
Let $L\in C(\overline{Q}\times U)$ be the \textit{running cost} and $\Psi\in C(I\times O)$ the \textit{terminal cost} defined as:

\begin{equation}\label{1-1-deftermcost}
    \Psi(t,x) = \begin{cases}
        g(t,x) & \text{if } (t,x)\in [t,t_1)\times O \\
        \psi(x) & \text{if } (t,x)\in \{t_1\}\times O  
    \end{cases}
\end{equation}

We define the \textit{payoff} $J$ as:

\begin{equation}\label{1-1-payoff1}
    J(t,x;u) = \int_t^{\tau}L(s,x(s),u(s)) \,ds + \Psi(\tau, x(\tau))
\end{equation}

Where $\tau$ is the exit time of $(s,x(s))$ from $\overline{Q}$, that is:

\begin{equation}\label{1-1-taudef}
    \tau = \begin{cases}
        \inf\{s\in [t,t_1)\,|\, x(s)\notin \overline{O}\} & \text{if } \exists s\in [t,t_1): x(s)\notin\overline{O} \\
        t_1 & \text{if }  x(s)\in\overline{O}\,\forall s\in [t,t_1)
    \end{cases}
\end{equation}

Then a control $u^{\ast}(\cdot)$ is \textit{optimal} if:

\begin{equation}\label{1-1-optimalcondition}
    J(t,x;u^{\ast}) \leq J(t,x;u) \quad \forall u\in L^{\infty}(I;U)
\end{equation}

Actually, we are being to generous with the control space. We have to impose a further condition on it, the \textit{switching condition}.
Let us assume that we have $u\in\mathcal{U}(t,x)$ and $u'\in\mathcal{U}(r,x(r))$ for $r\in[t,\tau]$. If we define:

\begin{equation}\label{1-1-switcond}
    \tilde{u}(s)=\begin{cases}
        u(s) & s\in[t,r) \\
        u'(s) & s\in[r,t_1]
    \end{cases}
\end{equation}

Then we impose:

\begin{equation}
    \tilde{u}_s\in\mathcal{U}(s,\tilde{x}(s)) \quad \forall s\in[t,\tilde{\tau}]
\end{equation}

Where $\tilde{x}$ is the solution to the control problem \ref*{1-1-syst1} with control $\tilde{u}$ and initial condition $x$, 
$\tilde{u}_s$ is the restriction of $\tilde{u}$ to $[s,t_1]$ and $\tilde{\tau}$ is the exit time of $(s,\tilde{x}(s))$ from $\overline{Q}$.
This condition assures that admissible controls can be replaced as the time evolves and the resulting control is still admissible. 

\section{Dynamic programming principle}

One way of tackling certain optimal control problems is via \textit{dynamic programming}. 
% The idea is to use a \textit{value function} as a tool to solve the control problem, finding sufficient and sometimes necessary conditions.
Let us define the \textit{value function}:

\begin{equation}\label{1-2-valuefunc}
    V(t,x)=\inf_{u\in \mathcal{U}(t,x)}J(t,x;u)
\end{equation}

For all $(t,x)\in\overline{Q}$. We get rid of the instance in which $V(t,x)=-\infty$ assuming $Q$ to be compact, or $L$ and $\Psi$ to be bounded below.
We aim at retrieving the argument which attains the infimum of \ref*{1-2-valuefunc}. In order to immerse this optimal control problem into 
a dynamic programming one we see the state of the system as the state of the variable and the control function as the decision function.
The basic idea behind dynamic programming techniques is to subdivide a problem into smaller problems, what does this mean in our context?
We will be able to find instantaneous the value function $V$ via a partial differential equation (PDE) called Hamilton-Jacobi-Bellman equation.

We start by stating and proving the following proposition, which provides us with an equivalent definition of the value function.

\begin{proposition}\label{1-2-propriformulation}
    For any $(t,x)\in\overline{Q}$ and any $r\in I$ then:

    \begin{equation}\label{1-2-riformulationvalue}
        V(t,x)=\inf_{u\in\mathcal{U}(t,x)}\left\{\int_t^{r\land\tau}L(s,x(s),u(s))\,ds+g(\tau,x(\tau))\chi_{\tau<r}+V(r,x(r))\chi_{r\leq\tau}\right\}
    \end{equation}

    \begin{proof}
        Value function less than rhs. If $r>\tau$ then $\tau<t_1$ and $\Psi(r\land\tau,x(r\land\tau))=g(\tau,x(\tau))$ and then \ref*{1-2-riformulationvalue}
        follows directly by definition. If $r\leq\tau$, let $\delta>0$ then there exists $u^1\in\mathcal{U}(r,x(r))$ such that:

        \[\int_r^{\tau^1}L(s,x^1(s),u^1(s))\,ds+\Psi(\tau^1,x^1(\tau^1))\leq V(r,x(r))+\delta\]

        Where $x^1$ is the state function corresponding to $u^1$ with initial condition $(r,x(r))$ and $\tau^1$ the first exit from $\overline{Q}$ of $(s,x^1(s))$.  
        By defining $\tilde{u}$ as for the switching condition \ref{1-1-switcond} we have $\tau^1=\tilde{\tau}$, because $\tau\geq r$ and then $\tilde{u}$ is $u^1$. Then:
        
        \begin{align*}
            V(t,x) & \leq V(t,x;\tilde{u}) \\
            & = \int_t^{\tilde{\tau}} L(s,\tilde{x}(s),\tilde{u}(s))\,ds+\Psi(\tilde{\tau},\tilde{x}(\tilde{\tau})) \\
            & = \int_t^{r} L(s,x(s),u(s))\,ds+\int_r^{\tau^1} L(s,x^1(s),u^1(s))\,ds+\Psi(\tau^1,x^1(\tau^1)) \\
            & \leq \int_t^{r} L(s,x(s),u(s))\,ds + V(r,x(r)) + \delta
        \end{align*}

    Since $\delta$ is arbitrary the first inequality is proved.
    
    \noindent Value function is bigger than rhs. Let $\delta>0$ and $U\in\mathcal{U}(t,x)$ such that:

    \[\int_r^{\tau}L(s,x(s),u(s))\,ds+\Psi(\tau,x(\tau))\leq V(t,x)+\delta\]

    Then: 

    \begin{align*}
        V(t,x) & \geq \int_r^{\tau}L(s,x(s),u(s))\,ds+\Psi(\tau,x(\tau)) - \delta \\
        % & = \int_t^{\tilde{\tau}} L(s,\tilde{x}(s),\tilde{u}(s))\,ds+\Psi(\tilde{\tau},\tilde{x}(\tilde{\tau})) \\
        & = \int_t^{r\land\tau} L(s,x(s),u(s))\,ds+\int_{r\land\tau}^{\tau} L(s,x(s),u(s))\,ds+\Psi(\tau,x(\tau)) -\delta\\
        % & \leq \int_t^{r} L(s,x(s),u(s))\,ds + V(r,x(r)) + \delta
        & = \int_t^{r\land\tau} L(s,x(s),u(s))\,ds+J(r,x(r))\chi_{r\leq\tau}+g(\tau,x(\tau))\chi_{\tau<r} - \delta\\
        & = \int_t^{r\land\tau} L(s,x(s),u(s))\,ds+V(r,x(r))\chi_{r\leq\tau}+g(\tau,x(\tau))\chi_{\tau<r} - \delta
    \end{align*}

    As $\delta$ is arbitrary we proved the proposition.
    \end{proof}
\end{proposition}


In the proof we used the concept of $\delta$-\textit{optimal control}, that is the control function $u\in\mathcal{U}(r,x(r))$ such that:

\[\int_r^{\tau^1}L(s,x^1(s),u^1(s))\,ds+\Psi(\tau^1,x^1(\tau^1))\leq V(r,x(r))+\delta.\]

This new representation allows us to find the so-called \textit{dynamic programming equation}. We have to impose that the value function 
is continuously differentiable, although this is not always the case. If differentiability fails, the notion of viscosity solution is needed. 

Let us first impose boundary conditions of the value function. Clearly if $t=t_1$ then:

\begin{equation}\label{1-2-boundcond1}
    V(t_1,x)=\psi(x)\,\forall x\in \overline{O}
\end{equation}

If $(t,x)\in [t_0,t_1)\times\partial O$ then the value function is $g$:

\begin{equation}\label{1-2-boundcond2}
    V(t, x) = g(t,x)
\end{equation}

Before stating the fundamental theorem which gives sufficient conditions 
for a solution to the optimal problem we follow a heuristic reasoning which will help our intuition. Under the hypothesis of continuous differentiability of the value function 
let us rewrite the dynamic programming principle as:

\begin{equation}
    \inf_{u\in\mathcal{U}}\left\{\frac{1}{h}\int_t^{(t+h)\land\tau} L(s,x(s),u(s))\,ds + \frac{1}{h}g(\tau,x(\tau))\chi_{\tau<t+h} + \frac{1}{h}\left[V(t+h,x(t+h))\chi_{\tau\geq t+h} - V(t,x)\right]\right\}=0
\end{equation}

Then if we formally let $h\to0$ we get:

\[\inf_{u\in\mathcal{U}}\left\{L(t,x(t),u(t)) + \partial_tV(t,x(t)) + D_xV(t,x(t))\cdot f(t,x(t),u(t))\right\}=0\]

Which can be rewritten as:

\begin{equation}\label{1-2-HJB1}
    -\frac{\partial}{\partial t}V(t,x) + H(t,x,D_xV(t,x))=0
\end{equation}

Where for $(t,x,p)\in \overline{Q}\times\R^n$ the Hamiltonian is defined as:

\begin{equation}\label{1-2-Hamiltonian1}
    H(t,x,p) = \sup_{v\in\R^n}\left\{-p\cdot f(t,x,v) - L(t,x,v)\right\}.
\end{equation}

Equation \ref{1-2-HJB1} turns out to be the main sufficient condition for the value function to be optimal.

Maybe only differentiability is needed (also for if, for only if we already know).

\begin{theorem}[Verification Theorem]\label{1-2-Verificationthe}
    Let $W\in C^1(\overline{Q})$ satisfy \ref{1-2-HJB1} and the boundary conditions \ref{1-2-boundcond1} and \ref{1-2-boundcond2} then:

    \[W(t,x)\leq V(t,x) \, \forall (t,x)\in\overline{Q}\]

    Moreover, there exists $u^{\ast}\in\mathcal{U}$ such that:

    \begin{equation}\label{1-2-uoptim}
        \begin{cases}
            L(s,x^{\ast}(s),u^{\ast}(s)) + f(s,x^{\ast},u^{\ast}(s))\cdot D_xW(s,x^{\ast}(s)) = - H(s,x^{\ast}(s),D_xW(s,x^{\ast}(s))) & \text{a.s. for } s\in[t,\tau^{\ast}] \\
            W(\tau^{\ast},x^{\ast}(\tau^{\ast})) = g(\tau^{\ast},x^{\ast}(\tau^{\ast})) & \text{if } \tau^{\ast}<t_1
        \end{cases}    
    \end{equation}

    if and only if $u^{\ast}$ is optimal and $W=V$.

    \begin{proof}
        Let $u\in\mathcal{U}$, then:

        \begin{align*}
            \Psi(\tau,x(\tau)) & = W(\tau,x(\tau)) = W(t,x(t)) + \int_t^{\tau}\frac{d}{ds}W(s,x(s)) \,ds \\
            & = W(t,x(t)) + \int_t^{\tau} \frac{\partial}{\partial t}W(s,x(s)) + \dot{x}(s)\cdot D_xW(s,x(s)) \,ds \\
            & = W(t,x(t)) + \int_t^{\tau} \frac{\partial}{\partial t}W(s,x(s)) + f(s,x(s),u(s))\cdot D_xW(s,x(s)) \,ds \\
            & \overset{\circledast}{\geq} W(t,x(t)) - \int_t^{\tau} L(s,x(s),u(s)) \,ds  
        \end{align*}

        Then:

        \[W(t,x(t)) \leq J(t,x;u)\]

        And therefore by taking the infimum over $\mathcal{U}$ and recalling $x(t)=x$ we get:

        \[W(t,x) \leq V(t,x)\]

        If furthermore $u^{\ast}$ satisfies \ref{1-2-uoptim} then the inequality $\overset{\circledast}{\geq}$ is an equality, and therefore:

        \[W(t,x) = J(t,x;u^{\ast})\]

        Which implies that $u^{\ast}$ is optimal and $W(t,x)=J(t,x;u^{\ast})=V(t,x)$. 
        The converse will be proved in a more general setting. In particular, only differentiability is needed.
        % Conversely, if $u^{\ast}$ is optimal:

        % \begin{align*}
        %     \color{red}{\int_t^{\tau} \frac{\partial}{\partial t}W(s,x^{\ast}(s)) + f(s,x^{\ast}(s),u^{\ast}(s)) \cdot D_xW(s,x^{\ast}(s)) \,ds & \geq - \int_t^{\tau^{\ast}} L(s,x^{\ast}(s),u^{\ast}(s)) \,ds  \\
        %     & \geq - \int_t^{\tau} L(s,x^{\ast}(s),u(s)) \,ds = }
        %     \color{black}{}
        % \end{align*}
    \end{proof} 
\end{theorem}


Theorem \ref*{1-2-Verificationthe} is an important tool in determining the explicit form of and optimal control. 
Indeed, condition \ref*{1-2-uoptim} can be restated as:

\begin{equation}\label{1-2-optimalityconditionu}
        u^{\ast}(s) \in \argmin_{v\in U} \left\{ f(s,x^{\ast}(s),v) \cdot D_xW(s,x^{\ast}(s)) + L(s,x^{\ast}(s),v)\right\}
\end{equation}

For almost all $s\in[t,t_1]$.

I have to prove verification theorem in more general case ($O\neq\R^n$).

We can express the optimality condition on $u$ \ref*{1-2-optimalityconditionu} in a differential inclusion form. 

\begin{corollary}
    A control $u^{\ast}$ is optimal if the corresponding state function $x^{\ast}$ satisfies:

    \begin{equation}
        x^{\ast} \in \left\{f(t,x,v) \,|\, v\in v^{\ast}(t,x)\right\}
    \end{equation}
\end{corollary}

\section{Pontryagin's principle and dynamic programming}

In the previous section we tackled the optimal control problem via dynamic programming.
As mentioned earlier this approach is of wide applicability and provides an implicit 
characterization of an optimal control. We now present another technique: the Pontryagin's principle.
As before, it will give rise to necessary condition on a control function 
to be optimal, but they'll come from a completely different perspective. We now present Pontryagin's 
principle in its full generality, and then we will see how it is connected with 
the dynamic programming approach. 

\subsection{Pontryagin's principle}

Pontryagin gives us an elegant and unintuitive way of solving problems of the kind:

\begin{equation}\label{1-3-problempon}
    \begin{cases}
        \dot{x}(s) = f(s,x(s),u(s)) & s\in [t,t_1] \\
        x(t) = x
    \end{cases}
\end{equation}

Where $u$ is bounded measurable into $U\subset\R^m$ and $O=\R^n$. Having defined the 
functional $J(t,x;u)$ as usual we define the \textit{control state Hamiltonian} 
as follows.

\begin{definition}
    The control state Hamiltonian of system \ref{1-3-problempon} is:

    \begin{equation}\label{1-3-hamiltonian}
        H(s,x,u,p) = - p \cdot f(s,x,u) - L(s,x,u)
    \end{equation}

    For all $s\in[t,t_1],\,x,p\in O,\,u\in U$.
\end{definition}

The variable $p$ is called \textit{costate} of the system. Pontryagin's principle 
gives us information about the costate under an optimal trajectory, which in turn will 
characterize the optimal control. 

\begin{theorem}\label{1-3-pontry}
    Let $u^{\ast}$ be an optimal control and $x^{\ast}$ its corresponding trajectory. 
    Then there exists a function $p^{\ast}:[t,t_1]\rightarrow O$ such that:

    \begin{equation}\label{1-3-pontry-x}
        \dot{x}^{\ast}(s) = D_p H(s,x^{\ast}(s),u^{\ast}(s),p^{\ast}(s)) 
    \end{equation}

    \begin{equation}\label{1-3-pontry-p}
        \dot{p}^{\ast}(s) =  - D_x H(s,x^{\ast}(s),u^{\ast}(s),p^{\ast}(s)) 
    \end{equation}

    And also:

    \begin{equation}\label{1-3-pontry-maxH}
        H(s,x^{\ast}(s),u^{\ast}(s),p^{\ast}(s)) = \sup_{v\in U} H(s,x^{\ast}(s),v,p^{\ast}(s))
    \end{equation}

    With:

    \begin{equation}\label{1-3-pontry-tras}
        p^{\ast}(t_1) = D \psi(x^{\ast}(t_1))
    \end{equation}
\end{theorem}

Thanks to this result we can determine an optimal control via the costate. By solving equation \ref{1-3-pontry-p} 
arisen in Theorem \ref{1-3-pontry} we can obtain the explicit form of $p^{\ast}$, and then 
retrieve $u^{\ast}$ from the maximization principle \ref{1-3-pontry-maxH}.


\subsection{Dynamic programming interplay}

As we just saw, Pontryagin's principle offers us an-even-though-quite-involved technique 
for finding an optimal control. As use of a control state Hamiltonian is crucial in this approach, it reminds 
us of the Hamiltonian defined in \ref{1-2-Hamiltonian1}. The similarity is also fortified by the maximization 
principle \ref{1-3-pontry-maxH}. We shall prove that this similarity unveils the direct link between 
Pontryagin's principle and the dynamic programming equation.

We will use the notion of differentiability of the value function. 
The classical notion of differentiability which we use is the following.

\begin{definition}
    V is differentiable in $(t,x)$ if there exists $V_t(t,x),V_x(t,x)\in\R$ such that:

    \begin{equation}
        \lim_{(h,k)\to(0,0)} \frac{1}{\abs{h}+\abs{k}}\abs{V(t+h,x+k)-V(t,x)-V_t(t,x)\cdot h-V_x(t,x)\cdot k}=0
    \end{equation}
\end{definition}

Differentiability is somewhat a strong hypothesis, but it allows us to prove the following
proposition. We must say that differentiability may easily fail in application, in such istances 
the notion of a weaker solution is needed, namely a viscosity solution.

\begin{theorem}
    Let $V$ be differentiable in $(t,x)\in Q$ and $u^{\ast}$ an optimal control
    such that $u^{\ast}\xrightarrow{s\to t}v$, then:

    \begin{equation}\label{1-3-dynamicprogrammeq}
        V_t(t,x) + L(t,x,v) + f(t,x,v)\cdot D_xV(t,x) = 0
    \end{equation}

    \begin{proof}
        Let $h>0$ s.t. $t+h<\tau$, then by \ref{1-2-propriformulation} we have:

        \[ V(t,x) = \int_t^{t+h} L(s,x(s),u^{\ast}(s)) \,ds + V(t+h,x(t+h))\]

        But because of differentiability we have:

        \[\lim_{h\to0}\frac{1}{\abs{h}}\abs{V(t+h,x(t+h))-V(t,x(t))} = V_t(t,x) + f(t,x,v)\cdot D_xV(t,x)\]

        Then we get:

        \begin{align}
        L(t,x,v) & = \lim_{h\to0} \frac{1}{\abs{h}} \int_t^{t+h} L(s,x(s),u(s)) \,ds = -\lim_{h\to0} \frac{1}{\abs{h}}\abs{V(t+h,x(t+h))-V(t,x(t))} \\
        & = - V_t(t,x) - f(t,x,v)\cdot D_xV(t,x)
        \end{align}
    \end{proof}
\end{theorem}

Furthermore, we impose existence and continuity of all derivatives of $f,L,g,\psi$.
The next theorem demonstrates that the costate in the Pontryagin Maximum
Principle is in fact the gradient in x of the value function v, taken along an optimal
trajectory. 

\begin{theorem}\label{1-3-pontrydyn}
    Let $u^{\ast}$ be an optimal right-continuous control and $x^{\ast}$ its corresponding trajectory. 
    Assume that the value function $V$ is differentiable at $(s,x^{\ast}(s))$ for $s\in[t,t_1)$. If we define:

    \begin{equation}\label{1-3-pontrydyn-defP}
        p(s) = D_x V(s,x^{\ast}(s)) 
    \end{equation}

    Then $p(s)$ satisfies \ref{1-3-pontry-p}, \ref{1-3-pontry-maxH} and \ref{1-3-pontry-tras}.

    \begin{proof}
        The trasversality condition \ref{1-3-pontry-tras} is straightforward from the definition of the  
        value function, which implies also the maximality condition \ref{1-3-pontry-maxH}. 
        We need to prove the "lagrangian multiplier condition" \ref{1-3-pontry-p}. 
        Let us drop all $^{\ast}$ and rewrite this differential equation:

        \begin{equation}\label{1-3-pontrydyn-rewritecond}
            \frac{d}{dt} p_j(s) = - \sum_{i=1}^n \frac{\partial}{\partial x_j}f_i(s,x(s),u(s))p_i(s) - \frac{\partial}{\partial x_j}L(s,x(s),u(s))
        \end{equation}

        This system admits solution $\overline{p}$ such that $\overline{p}(s)=D_xV(s,x(s))$; let us show that $\overline{p}(s)=p(s)$. Let $u_s$ be 
        the restriction of $u$ to $[r,t_1)$, which is admissible by assumption. We have:

        \[V(s,y) \leq J(s,y;u_s) \,\forall y\in\R^n\]

        Then, because $u$ is optimal, $y\mapsto J(s,y;u_s) - V(s,y)$ has its global minimum in $y=x(s)$, which implies by differentiability:

        \begin{equation}\label{1-3-pontrydyn-derconV}
            D_x V(s,x(s)) = D_x J(s,x(s);u_s)
        \end{equation}

        Then we prove that $\overline{p}(s) = D_x J(s,x(s);u_s)$. We denote $x(r)$ the solution starting at $x(s)$ at time $r$. Because $L\in C^1$ then for all $i=1,\dots,n$:

        \begin{equation}
            \frac{\partial}{\partial x_i}J(s,x(s),u) = \sum_{j=1}^n \int_s^{t_1} \left(L_{x_j}(r,x(r),u(r))\frac{\partial x_j(r)}{\partial x_i}\right) \,dr + \psi_{x_j}(x(t_1))\frac{\partial x_j(t_1)}{\partial x_i}
        \end{equation} 

        But then:

        \begin{align}
            \overline{p}_i(s)  & = \sum_{j=1}^n \frac{\partial x_j(s)}{\partial x_i}\overline{p}_j(s) = \sum_{j=1}^n \frac{\partial x_j(t_1)}{\partial x_i}\overline{p}_j(t_1) - \int_s^{t_1} \frac{d}{dr}\left(\sum_{j=1}^n \frac{\partial x_j(r)}{\partial x_i}\overline{p}_j(r)\right) \,dr \\
            & = \sum_{j=1}^n \frac{\partial x_j(t_1)}{\partial x_i}\psi(x(t_1)) - \sum_{j=1}^n \int_s^{t_1} \frac{d}{dr}\left(\frac{\partial x_j(r)}{\partial x_i}\right)\overline{p}_j(r) + \frac{\partial x_j(r)}{\partial x_i}\frac{d}{dr}\left(\overline{p}_j(r)\right) \,dr 
        \end{align}

        But under the integral:

        \[\int_s^{t_1} \frac{d}{dr}\sum_{l=1}^n \frac{\partial x_j(r)}{\partial x_i}\left(\frac{\partial}{\partial x_j}f_l(r,x(r),u(r))\overline{p}_j(r) - \sum_{l=1}^n \frac{\partial x_j(r)}{\partial x_i}\frac{\partial}{\partial x_j}f_l(r,x(r),u(r))\overline{p}_j(r) - \frac{\partial}{\partial x_j}L(r,x(r),u(r))\right) \,dr\]

        Then we get:

        \begin{align}
            \overline{p}_i(s) & = \sum_{j=1}^n \left(\frac{\partial x_j(t_1)}{\partial x_i}\psi(x(t_1)) + \int_s^{t_1} \frac{\partial}{\partial x_j}L(r,x(r),u(r))\frac{\partial x_j(r)}{\partial x_i} \,dr\right) \\
            & = \frac{\partial}{\partial x_i}J(s,x(s);u_r)
        \end{align}
    \end{proof}
\end{theorem}



\section{Existence}

We now prove an existence theorem for optimal controls.
We study the fixed time interval case with $O=\R^n$ and the function $f$ linear in $v$.
Furthermore, we impose convexity of $L$ in $v$. Under these assumptions a classical variational
argument proves the optimal control existence.
\begin{theorem}
    Let $U$ compact and convex, $f_1,f_2\in C^1(\overline{Q}\times U)$ such that $f(t,x,v)=f_1(t,x)+f_2(t,x)v$ and $\partial_xf_{1},\partial_xf_2,f_2$ bounded.
    Let also $L\in C^1(\overline{Q}\times U)$ and $L(t,x,\cdot)$ be convex for all $(t,x)\in\overline{Q}$ and the temrinal cost $\phi\in C(\R^n)$. Then there exist an optimal control $u^{\ast}(\cdot)$.

    \begin{proof}
        Let $u_n$ a minimizing sequence such that:

        \begin{equation}
            \lim_{n\to+\infty} J(t,x;u_n) = V(t,x)
        \end{equation}

        Let $x_n(\cdot)$ be the solutions to \ref{1-1-syst1} with $u=u_n$. If we show both sequence to converge respectively 
        (weakly) to $u^{\ast}$ and uniformly $x^{\ast}$ (along subsequences) such that the latter is again the solution to \ref{1-1-syst1} with $u=u^{\ast}$, then:

        \[J(t,x;u_n) = \int_t^{t_1} L(s,x^{\ast}(s),u_n(s)) \,ds + \int_t^{t_1} L(s,x_n(s),u_n(s)) - L(s,x^{\ast}(s),u_n(s)) \,ds + \phi(x_n(t_1))\]

        But then:

        \[\liminf_{n\to+\infty} \int_t^{t_1} L(s,x_n(s),u_n(s)) - L(s,x^{\ast}(s),u_n(s)) \,ds = 0\]

        And $\psi(x_n(t_1))\xrightarrow{n\to+\infty}\psi(x^{\ast}(t_1))$. But then:

        \begin{equation}\label{1-4-existenceliminf}
            \liminf_{n\to+\infty} J(t,x;u_n) = \liminf_{n\to+\infty} \int_t^{t_1} L(s,x^{\ast}(s),u_n(s)) \,ds \geq \int_t^{t_1} L(s,x^{\ast}(s),u^{\ast}(s)) \,ds 
        \end{equation}

        Because $L$ is convex in $u$\footnote{Explained in remark \ref{1-4-reamrkconvexity}}. Therefore:

        \[V(t,x)\leq J(t,x;u^{\ast}) \leq \liminf_{n\to+\infty} J(t,x;u_n) = V(t,x)\]

        We need to prove convergence of $x_n$ and $u_n$. Because $U$ is compact and \textit{convex} then $L^{\infty}([t,t_1];U)$ is 
        weakly sequentially compact. For what concerns $x_n$ we use Ascoli-Arzela's theorem to show that is admits a uniformly convergent subsequence. Being 
        uniformly limitated comes from:

        \begin{align}
            \abs{x_n(s)} & \leq \abs{x_n(s)-x}+\abs{x} = \abs{x} + \int_t^s \frac{d}{dr}\abs{x_n(r)} \,dr \\
            & = \abs{x} + \int_t^s \abs{f_1(r,x_n(r)) + f_2(r,x_n(r))\cdot u_n(r)}\,dr \\
            & \color{red}{\leq} \abs{x} + \int_t^s \norm{\partial_xf_1}_{\infty}\abs{x_n(r)} + \norm{f_2}_{\infty}\abs{u_n} \,dr \\
            & \leq C + K\left( \int_t^s \abs{x_n(r)} \,dr \right)   
        \end{align}

        Then by Gronwall's lemma $x_n$ is uniformly limitated. Equicontinuity comes from the uniform buondedness of the
        derivative $\dot{x}_n(s)$. Therefore, we know that there exist the weak limit $u^{\ast}$ and the uniform
        limit $x^{\ast}$. The latter is the solution of \ref{1-1-syst1} with $u=u^{\ast}$. Indeed:

        \begin{align*}
                x_n(s) & = x + \int_t^s\frac{d}{dr}x_n(r) \,dr\\
                & = x + \underbrace{\int_t^sf_1(r,x_n(r)) + f_2(r,x_n(r))u^{\ast}(r) \,dr}_{A_n} + \underbrace{\int_t^s \left[f_2(r,x_n(r)) - f_2(r,x^{\ast}(r))\right]\left[u_n(r)-u^{\ast}(r)\right] \,dr}_{B_n} \\
                & + \underbrace{\int_t^s f_2(r,x^{\ast}(r))\left[u_n(r)-u^{\ast}(r)\right] \,dr}_{C_n}
        \end{align*}
        
        Letting $n\to+\infty$ we get $B_n$ (by weak convergence and boundedness of $f_2$) and $C_n$ (by weak convergence) going to $0$ 
        while we obtain:

        \[A_n\xrightarrow{n\to+\infty}\int_t^sf_1(r,x^{\ast}(r))+f_2(r,x^{\ast}(r))u^{\ast}(r)\,dr\]
        
        And therefore the thesis.
    \end{proof}
\end{theorem}


\begin{remark}\label{1-4-reamrkconvexity}
    In the proof we asserted inequality \ref{1-4-existenceliminf} by convexity of the running cost. Indeed, 
    by convexity and being $C^1$:

    \[L(s,x(s),u_n(s)) \geq L(s,x(s),u^{\ast}(s)) + \left[u_n(s) - u^{\ast}(s)\right]L_u(s,x(s),u^{\ast}(s))\]

    Then by integrating and taking $\liminf_{n\to+\infty}$ we get the inequality (using weak convergence of $u_n$). 
\end{remark}

An existence result can also be proved in the context of $O\neq\R^n$, where the cost function
to be minimized has the form:

\begin{equation}
    J(t,x;u)=\int_t^{\tau} L(s,x(s),u(s)) \,ds + \Psi(\tau,x(\tau))
\end{equation}

with dynamical system described by $f(s,x,u)=u$. Therefore, the value function is:

\begin{equation}
    V(t,x) = \inf_{u\in\mathcal{U}} \left\{\int_t^{\tau} L(s,x(s),\dot{x}(s)) \,ds + \Psi(\tau,x(\tau))\right\}
\end{equation}

Under technical assumptions an optimal control exists, and the proof follows
the same procedure as before: we take a minimizing sequence, show its (weak) convergence
and its corresponding state evolution (uniform) convergence, then we show 
optimality, slightly modifying the proof because of convergence issues.

\section{Infinite horizon}

A particularly interesting version of the maximization problem arising from optimal 
control is with infinite horizon. Let us consider the usual Cauchy' problem:

\begin{equation}
    \begin{cases}
        \dot{x}(s) = f(s,x(s),u(s)) & s\in[t,t_1] \\
        x(t) = x
    \end{cases}
\end{equation}

If we set $t_1=+\infty$ we get an infinite horizon problem. Thus, the maximization 
problem becomes:

\begin{equation}
    \inf_{u\in\mathcal{U}} \int_t^{\tau} L(s,x(s),u(s)) \, ds + g(\tau,x(\tau))\chi_{\tau<+\infty}
\end{equation}

Where $\tau$ is the exit time of $x(\cdot)$ from $\overline{O}$.\footnote{More precisely, is the exit time of $(s,x(s))$ from $\overline{Q}$.}

\section{Proof of Pontryagin's principle}

We will show the maximum Pontryagin's principle in the simple context of no running cost. 
A cleaver reconstruction of the non-zero running cost problem as a zero running cost one 
will enlarge the thesis to this situation. The first concept we will need is the variation 
of a control.

\begin{definition}
    Given $u\in\mathcal{U}$. For $\epsilon,r>0$ such that $0<r-\epsilon<r$ and $a\in U$ 
    we define the \textit{simple variation} $u_{\epsilon}\in\mathcal{U}$ such that:

    \begin{equation}\label{1-proofpontry-defvar}
        u_{\epsilon}(t) = \begin{cases}
            a & s\in(r-\epsilon,r) \\
            u(s) & s\notin(r-\epsilon,r)
        \end{cases}
    \end{equation}
\end{definition}

By defining the matrix $A:[0,+\infty)\rightarrow\R^{n\times n}: s\mapsto D_xf(s,x(s),u(s))$ we state 
the following lemma.

\begin{lemma}
    Let $x_{\epsilon}$ be solution of:

    \begin{equation}\label{1-proofpontry-dynamprob}
        \begin{cases}
            \dot{x}_{\epsilon}(s) = f(s,x_{\epsilon}(s),u_{\epsilon}(s)) & s\in[t,t_1] \\
            x_{\epsilon}(t) = x
        \end{cases}
    \end{equation}

    Then the solution is:

    \begin{equation}\label{1-proofpontry-solvar}
        x_{\epsilon}(s) = x(s) + \epsilon y(s) + o(\epsilon)\, \epsilon\to0
    \end{equation}

    Where $y\equiv0$ on $[t,r]$ and:

    \begin{equation}
        \begin{cases}
            \dot{y}(s) = A(s)y(s) & s\in[r,t_1]\\
            y(r) = y^r
        \end{cases}
    \end{equation}

    With $y^r=f(x(r),a)-f(x(r),u(r))$.

    \begin{proof}
        Let us divide the proof in three cases.

        \begin{itemize}
            \item $s\in[t,r-\epsilon]$: then $y(t)=0$ and $u_{\epsilon}(t)=u(t)$, therefore:
            
            \[x_{\epsilon}(t)=x(t)=x(t)+\epsilon y(t)+o(\epsilon)\]

            \item $s\in(r-\epsilon,r)$: then we have:
            
            \[x_{\epsilon}(s)-x(s) = \int_{r-\epsilon}^s f(w,x_{\epsilon}(w),u_{\epsilon}(w))-f(w,x(w),u(w)) \,dw + o(\epsilon)\]

            Which is a little \textit{o} of $\epsilon$ (because $f$ is continuous). 

            \item $s\in[r,t_1]$: from before if $s=r$ then:
            
            \begin{align}
                x_{\epsilon}(r)-x(r) & = \int_{r-\epsilon}^r f(w,x_{\epsilon}(w),u_{\epsilon}(w))-f(w,x(w),u(w)) \,dw + o(\epsilon) \\
                & = \lim_{w\to s}[f(w,x_{\epsilon}(w),u_{\epsilon}(w))-f(w,x(w),u(w))]\epsilon+ o(\epsilon) \\
                & = y^s\epsilon + o(\epsilon)
            \end{align}

            \color{red}{If $s>r$ then:}


        \end{itemize}
    \end{proof}
\end{lemma}

Let us now prove Pontryagin's principle with no running cost. The payoff functional is:

\begin{equation}\label{1-proofpontry-norunfunct}
    J(t,x;u) = \psi(x(t_1))
\end{equation}

And therefore the Hamiltonian is:

\begin{equation}\label{1-proofpontry-norunham}
    H(s,x,u,p) = - f(s,x,u)\cdot p 
\end{equation}

\begin{theorem}\label{1-proofpontry-theo}
    There exists a function $p^{\ast}:[t,t_1]\rightarrow \R^n$ such that:

    \begin{equation}
        \dot{p}^{\ast}(s) = -D_x H(s,x^{\ast}(s),u^{\ast}(s),p^{\ast}(s))\, s\in[t,t_1]
    \end{equation}

    together with the maximization:

    \begin{equation}
        H(s,x^{\ast}(s),u^{\ast}(s),p^{\ast}(s)) = \sup_{v\in U} H(s,x^{\ast}(s),v,p^{\ast}(s))
    \end{equation}

    and the trasversality condition:

    \begin{equation}
        p^{\ast}(t_1) = D\psi(x(t_1))
    \end{equation}

    \begin{proof}
        Let us drop all the $^{\ast}$. Let $p$ be the unique solution of:

        \begin{equation}
            \begin{cases}
                \dot{p}(s) = - A'(s)\cdot p(s) & s\in[t,t_1] \\
                p(t_1) = D\psi(x(t_1))
            \end{cases}
        \end{equation}

        It exists and is unique because the latter is a linear differential equation with integrable coefficient. 
        We already satisfy the trasversality condition and the adjoint dynamics. We prove the 
        maximization principle. Let $a\in U$. We define the variation $u_{\epsilon}$ for $\epsilon,r\in(t,t_1)$ as before.
        Since $\epsilon\mapsto  J(t,x;u_{\epsilon})$ for $\epsilon\in[0,1]$ has a maximum in $\epsilon=0$ we have:
        
        \begin{equation}\label{1-proofpontry-dervar0}
            \frac{d}{d\epsilon}J(t,x;u_{\epsilon})\leq0
        \end{equation}

        Computing the derivative, using \ref{1-proofpontry-defvar}:

        \begin{align}
            \frac{d}{d\epsilon}J(t,x;u_{\epsilon})\big|_{\epsilon=0} & = \frac{d}{d\epsilon}\psi(x_{\epsilon}(t_1))\big|_{\epsilon=0} \\
            & = \frac{d}{d\epsilon}\psi(x(t_1) + \epsilon y(t_1) + o(\epsilon)) = D\psi(x(t_1))\cdot y(t_1) \\
            & = p(t_1)\cdot y(t_1) = p(r)\cdot[f(r,x(r),a)-f(r,x(r),u(r))]
        \end{align}

        Where the last equality comes from:

        \begin{align*}
            \frac{d}{ds}\left(p(s)\cdot y(s)\right) & = \dot{p}(s)\cdot y(s) + p(s)\cdot \dot{y}(s) \\
            & = -A'(s)\cdot p(s)\cdot y(s) + p(s) \cdot A(s)\cdot y(s) = 0
        \end{align*}

        Therefore, by plugging into \ref{1-proofpontrydervar0} we get:

        \[0\geq p(r)\cdot[f(r,x(r),a)-f(r,x(r),u(r))]\]

        Which implies:

        \[H(r,x(r),a,p(r)) = f(r,x(r),a) \cdot p(r) \leq f(r,x(r),u(r)) \cdot p(r) = H(r,x(r),u(r),p(r))\]
    \end{proof}
\end{theorem}

Given Pontryagin's principle for no running cost problems we can extend the result to the general case:

\begin{equation}
    J(t,x;u) = \int_t^{t_1} L(s,x(s),u(s)) \,ds + \psi(x(t_1))
\end{equation}

Where the Hamiltonian is:

\begin{equation}
    H(s,x,u,p) = f(s,x,u)\cdot p + L(s,x,u)
\end{equation}

Indeed, theorem \ref{1-proofpontry-theo} holds also under these conditions. We rewrite the problem as it has no 
running cost and then apply the theorem. Let us define $x^{n+1}$ as:

% \begin{equation}
%     \begin{cases}
%         \dot{x}^{n+1}(s) = L(s,x^{n+1}(s),u(s)) & s\in [t,t_1] \\
%         x^{n+1}(0) = 0 
%     \end{cases}
% \end{equation}

\begin{equation}
    x^{n+1}(s) = \int_t^s L(w,x(w),u(w)) \,dw
\end{equation}

Then by defining $\overline{f}, \overline{g}, \overline{x},\overline{x}(s)$ as:

\begin{equation}
    \overline{f}(s,x,u)=\begin{bmatrix}
        f(s,x,u) \\
        L(s,x,u)
    \end{bmatrix},\,\overline{x}=\begin{bmatrix}
        x \\
        0
    \end{bmatrix},\,\overline{x}(s)=\begin{bmatrix}
        x(s) \\
        x^{n+1}(s)
    \end{bmatrix},\,\overline{g}(\overline{x}(t_1)) = g(x(t_1)) + x^{n+1}(t_1)
\end{equation}

Thus, the problem has no running cost. We can apply the theorem and, noticing that 
$p^{n+1}\equiv1$ we get the thesis.

\newpage

\chapter{Stochastic Optimal Control}

We now study stochastic optimal control. The system we aim to control is governed by stochastic differential equations. 

Short description of what is done in this chapter.

\section{Markov diffusion process}

I now recall some definitions, give new ones and set the notation. Let $\Sigma\subseteq\R^n$ and $\mathcal{B}(\Sigma)$ 
the associated Borel $\sigma$-algebra. Let $(\Omega, \mathcal{F}, P)$ a general probability space. 
Given $x(s,\omega)$ a $\Sigma$-valued random process from $I_0=[t_0,t_1)$ and $(\Omega,\mathcal{F})$, let us denote by:

\[P(C\,|\,x(s_1),\dots,x(s_m)),\,C\in\mathcal{F}\]

The conditional probability of $C$ given the sigma algebra $\bigvee_{i=1}^m\sigma(x(s_i))$.

\begin{definition}\label{2-1-markovprocessdef}
    A stochastic process $x$ satisfies the Markov property if there exists a function 
    $p:I_0\times \Sigma\times I_0\times  \mathcal{B}(\Sigma)\rightarrow \R$ such that:

    \begin{enumerate}
        \item For all $t,s,B$ the function $x\mapsto p(t,x,s,B)$ is borel measurable on $\Sigma$
        \item For all $t,x,s$ the function $A\mapsto p(t,x,s,B)$ is a probability measure on $(\Omega,\mathcal{F})$
        \item The Chapman-Kolmogorov equation holds for all $s,t,r\in I_0$ such that $t<r<s$:
        \begin{equation}\label{2-1-markovprocessdef-chapkol}
            p(t,x,s,B) = \int_{\Sigma} p(r,y,s,B)\,p(t,x,r,dy)
        \end{equation}
    \end{enumerate}

    And such that for all $r,s\in I_0$ where $r,s$ and for all $B\in\mathcal{B}(\Sigma)$ then:

    \begin{equation}\label{2-1-markovprocessdef-condonp}
        P(x(s)\in B\,|\,\mathcal{F}_r^x) = p(r,x(r),s,B)
    \end{equation}

    Where $\mathcal{F}_r^x=\sigma\left(x(l)\,:\,l\in[t_0,r]\right)$.
\end{definition}

Function $p$ is called \textit{Markov Transition Kernel}. We shall see a Markov transition kernel as 
the probability that the system starting from $x$ at time $t$ will be in $B$ at time $s$. This heuristic interpretation 
clarifies the following notation:

\begin{equation}
    E_{tx}\phi(x(s)) = \int_{\Sigma} \phi(y)\,p(t,x,s,dy) 
    % p(t,x,s,B) = P(x(s)\in B\,|\,)
\end{equation}

For a real valued borel-measurable function $\phi$. Given a Markov process $x$ we can define a family of 
linear operators associated to it. Let $t<s$, hereafter all time indices will always be in $I_0$, and define:

\begin{equation}
    T_{t,s}\phi(x) = \int_{\Sigma} \phi(y)\,p(t,x,s,dy) = E_{tx}\phi(x(s))
\end{equation}

Integrability assumptions on $\phi$ vary from case to case. For now, we can take $\phi$ to be bounded. 
Because of Chapman-Kolmogorov equation \ref{2-1-markovprocessdef-chapkol} the family $(T_{t,s})_{t,s\in I_0}$ 
satisfies the property:

\begin{equation}\label{2-1-propT}
    T_{tr}\left[T_{rs}\phi\right] = T_{ts}\phi
\end{equation}

For all $t<r<s$. This family of linear operators defines another operator, the \textit{backward evolution operator}.
Let $A:\left\{\Phi:I_0\times\Sigma\rightarrow\R\right\}\rightarrow \R$:

\begin{equation}\label{2-1-backwarddef}
    A\Phi(t,x) = \lim_{h\to0+} \frac{E_{tx}\Phi(t+h,x(t+h))-\Phi(t,x)}{h}
\end{equation}

provided that the limit exists. We define $\mathcal{D}(A)$ the space of functions such that limit \ref{2-1-backwarddef} exists. 
The following holds.

\begin{proposition}
    Let $A$ as before, then for all $\Phi\in\mathcal{A}$ the following hold:

    \begin{enumerate}
        \item $\Phi,\frac{\partial\Phi}{\partial t}$ and $A\Phi$ are continuous
        \item For all $t,s\in\overline{I}_0$, $t<s$ then:
        
        \[E_{tx}\abs{\Phi(s,x(s))}< +\infty,E_{tx}\int_t^s\abs{A\Phi(r,x(r))}\,dr < +\infty\]
        
        \item Dynkin's formula holds for all $t<s$:
        
        \begin{equation}\label{2-1-dynkform}
            E_{tx}\Phi(s,x(s)) - \Phi(t,x) = E_{tx}\int_t^s A\Phi(r,x(r)) \,dr
        \end{equation}
    \end{enumerate}

    % \begin{proof}
    %     I prove Dynkin's formula in the case of $T_{ts}$ being a Feller semigroup. 
    %     %Dynkin, E. B., Markov processes,  

    % \end{proof}
\end{proposition}

Dynkin's formula can be proved in different instances, subject to the nature of the random process. We will see 
that it is a natural consequence of Ito formula for continuous state space processes.

If the random process $x$ is autonomous (time-homogeneous) then the linear operator family is a 
semigroup. Recall that a Markov process is homogeneous if for all $t<s$ in $I_0$ then:

\[p(t,x,s,B) = p(0,x,s-t,B)\]

If so, by calling $T_s=T_{0s}$ property \ref{2-1-propT} is:

\begin{align}
    T_{s+r}\phi(x) & = \int_{\Sigma} \phi(y)\,p(0,x,s+r,dy) \\
    % & = \int_{\Sigma} \phi(y)\,p(r,x,s,dy) \\
    & = \int_{\Sigma}\phi(y)\int_{\Sigma} p(r,z,r+s,dy)\,p(0,x,r,dz) \\
    & = \int_{\Sigma}\int_{\Sigma} \phi(y) p(r,z,r+s,dy)\,p(0,x,r,dz) \\
    & = \int_{\Sigma}\int_{\Sigma} \phi(y) p(0,z,s,dy)\,p(0,x,r,dz) \\
    & = \int_{\Sigma} T_{s}\phi(z) \,p(0,x,r,dz) \\
    & = T_r\left[T_s\phi(x)\right].
\end{align}

While the backward evolution operator analogous is called the \textit{generator} and is defined as:

\begin{equation}
    G\phi(x) = - \lim_{h\to0^+} \frac{T_h\phi(x) - \phi(x)}{h}
\end{equation}

With $D(G)$ as $\mathcal{D}(A)$ before. It is worth noting that, formally, the following equality holds:

\begin{equation}
    A\Phi = \frac{\partial \Phi}{\partial t} - G\Phi(t,\cdot)
\end{equation}

This relation links the two operators and the autonomous to the non-autonomous case. We now turn our attention to 
a subset of Markov processes: diffusion processes. A diffusion process is a Markov process whose paths are continuous. 
More formally. 

\begin{definition}
    A diffusion process $x:\overline{I}_0\times\Omega\rightarrow\Sigma$ is a almost surely continuous Markov process with Markov transition kernel $p$ such that:

    \begin{itemize}
        \item For every $\epsilon>0$:
        
        \begin{equation}
            \lim_{h\to0^+} \int_{\abs{x-y}>\epsilon} \,p(t,x,t+h,dy) = 0
        \end{equation}

        \item There exist functions $a_{ij}(t,x),f_{ij}(t,x)$ for $(t,x)\in\overline{Q}_0$ and $i,j=1,\dots,n$ such that for every $\epsilon>0$:
        
        \begin{equation}
            \lim_{h\to0^+} \int_{\abs{x-y}\leq\epsilon} (y_i-x_i)\,p(t,x,t+h,dy) = f_i(t,x)
        \end{equation}

        And:

        \begin{equation}
            \lim_{h\to0^+} \int_{\abs{x-y}\leq\epsilon} (y_i-x_i)(y_j-x_j)\,p(t,x,t+h,dy) = a_{ij}(t,x).
        \end{equation}
    \end{itemize}

    These limits are intended uniformly.
\end{definition}

Functions $f=(f_1,\dots,f_n)$ and $a=(a_{ij})_{ij}$ are respectively called local drift and local covariance coefficients.

How does the backward evolution operator, and the generator in the autonomous case, adapt to this situation? 
To answer this question we reduce our problem to a stochastic differential one by relying on the differential structure of a diffusion process. 
Give the local drift and covariance $f,a$ of a diffusion process $x$ we claim that it satisfies:

\begin{equation}\label{2-1-SDEdef}
    dx(s) = f(s,x(s))ds + \sqrt{a}(s,x(s))dw(s) 
\end{equation}

Clearly we have to impose further conditions of the stochastic differential equation's coefficients to ensure existence of a solution. 
In particular, we want those coefficients to be Lipshitz and sub-linearly growing with respect to the second variable. In equation \ref{2-1-SDEdef} 
We define the square root of $a$ as a function $\sqrt{a}=\sigma$ such that:

\begin{equation}
    \sigma(t,x)\cdot\sigma'(t,x) = a(t,x)
\end{equation}

We recall that under existence hypothesis for every $\Phi\in C^{1,2}(\overline{Q}_0)$ Ito's formula holds:

\begin{equation}\label{2-1-itotox}
    d\Phi(s,x(s)) = \Phi_s(s,x(s))ds + \sum_{i=1}^n \Phi_{x^i}(s,x(s)) dx^i(s) + \frac{1}{2}\sum_{i,j=1}^n \Phi_{x^ix^j}(s,x(s)) d[x,x]^{ij}(s)
\end{equation}

where $[x,y]$ is the covariation of process $x$ and $y$. Recall that this relation has always to be intended in integral form, that is:

\begin{equation}
    \Phi(s,x(s)) = \Phi(t,x) + \int_t^s\Phi_s(r,x(r))\,dr + \sum_{i=1}^n \int_t^s \Phi_{x^i}(r,x(r))\,dx^i(r) + \frac{1}{2}\sum_{i,j=1}^n \int_t^s\Phi_{x^ix^j}(r,x(r))\,d[x,x]^{ij}(r).
\end{equation}

Via this relation we can reconstruct Dynkin's formula in this setting. By defining the operator $A$ as in \ref{2-1-dynkform} we have:

\begin{align}
    \Phi(s,x(s)) & = \Phi(t,x) + \int_t^s\Phi_s(r,x(r))\,dr \\
    & + \sum_{i=1}^n \left[\int_t^s \Phi_{x^i}(r,x(r))f_i(r,x(r))\,dr + \sum_{j=1}^n\int_t^s\Phi_{x^i}(r,x(r))\sigma_{ij}(r,x(r))\,dw^j(r)\right] \\
    & + \frac{1}{2}\sum_{i,j=1}^n \sum_{l=1}^n\int_t^s\Phi_{x^ix^j}(r,x(r))\sigma_{il}(r,x(r))\sigma_{jl}(r,x(r))\,dr \\
    & = \Phi(t,x) + \int_t^s \Phi_s(r,x(r)) + D_x\Phi\cdot f(r,x(r)) + \frac{1}{2}D^2_x\Phi\cdot a(r,x(r)) \,dr \\
    & + \int_t^s D_x\Phi\cdot\sigma(r,x(r)) \,dw(r) \\
    % & = \Phi(t,x) + \int_t^s A\Phi(r,x(r))\,dr + \int_t^s D_x\Phi\cdot\sigma(r,x(r)) \,dw(r)
\end{align}

but the last (stochastic) integral can be seen as a martingale. In particular, if we take $\Phi$ to have polynomial growth of some order $m$:

\begin{equation}
    \abs{\Phi(t,x)} \leq K(1+\abs{x}^m)\,\forall(t,x)\in\overline{Q}_0
\end{equation}

then $D_x\Phi\cdot\sigma\in\mathcal{L}^2(I_0)$, where:

\[\mathcal{L}(I) = \left\{x:I\times\Omega\rightarrow\Sigma\,|\,E\int_Ix(s)\,ds<\infty\right\}\]

and therefore its stochastic integral is a martingale (with respect to the canonical filtration associated to the Brownian motion $w$). 
Therefore, if we take the (conditional) expectation:

\begin{equation}
    E_{tx}\Phi(s,x(s)) = \Phi(t,x) + E_{tx}\int_t^s \Phi_s(r,x(r)) + D_x\Phi\cdot f(r,x(r)) + \frac{1}{2}D^2_x\Phi\cdot a(r,x(r))\,dr.
\end{equation}

It is now natural to define the operator $A:C_p^{1,2}(\overline{Q}_0)\rightarrow\R$ as:

\begin{equation}\label{2-1-newdefA}
    A\Phi(r,x(r)) = \Phi_s(r,x(r)) + D_x\Phi\cdot f(r,x(r)) + \frac{1}{2}D^2_x\Phi\cdot a(r,x(r))
\end{equation}

where $C_p^{1,2}(I)$ is the family of functions $g$ from $I$ into $\R$ such that $g,g_s,g_{x_i},g_{x_ix_j}$ are continuous and with polynomial growth.

\begin{remark}
    Be careful that the stochastic integral:

    \[\int_t^s D_x\Phi\cdot\sigma(r,x(r))\,dw(r)\]

    is a martingale because $x$ satisfies:

    \[E_{tx}\abs{x(r)}^m\leq C_m(1+\abs{x}^m)\,\forall r\in I_0\]

    as it is solution of the SDE \ref{2-1-SDEdef}.\footnote{This is a standard result in SDE theory.}
\end{remark}

Consequently the generator $G$ of the time-homogeneous case is defined as:

\begin{equation}
    G\Phi(x) = -\frac{1}{2}\sum_{i,j=1}^n a_{ij}(x)\Phi_{x_ix_j}(x) - \sum_{i=1}^n f_i(x)\Phi_{x_i}(x)
\end{equation}


THEN CONTROLLED MARKOV PROCESSES.

I HAVE TO SEARCH FOR A MORE GENERAL VERSION OF EXISTENCE THEOREM for markov control policies.

THEN DYNAMIC PROGRAMMING APPROACH AND VERIFICATION THEOREM.

THEN EXAMPLE IF TIME ALLOWS.

THEN REREAD, REPEAT. I DONT THINK I'LL DO SLIDES. MAYBE THEY CAN BE A GOOD ASSET. LESS INFO, LESS THINGS TO KNOW PERFECTLY. BUT I HAVE TO KNOW EVERYTHING!!!\

\section{Markov control processes}

So far we talked about Markov processes without specifying any kind of control. A control process in any 
stochastic process $u:\Omega\rightarrow U$, where $U$ is the control space, that influences the evolution of the random process $x$. 
Formally, let $Q=I_0\times O$ and $u$ as before and define:

\begin{equation}\label{2-1-controlledSDE}
    \begin{cases}
        dx(r) = f(r,x(r),u(r))dr + \sigma(r,x(r),u(r))dw(r)\,r\in I_0 \\
        x(t) = x
    \end{cases}
\end{equation}

where $U\subset\R^m$ closed, $f,\sigma\in C(\overline{Q}_0\times U)$, $f(\cdot,\cdot,v),\sigma(\cdot,\cdot,v)$ belong to $C^1(\overline{Q}_0)$ for all $v\in U$, 
such that there exists $C>0$ such that:

\begin{align}
        \abs{f_t}+\abs{f_x} \leq C,\,\abs{\sigma_t}+\abs{\sigma_x}\leq C \\
        \abs{f(t,x,v)}\leq C(1+\abs{x}+\abs{v}) \\
        \abs{\sigma(t,x,v)}\leq C(1+\abs{x}+\abs{v})
\end{align}


We can relax the assumption by imposing Lipschitz condition on $t$ and $x$ for every fixed $v$. Furthermore,
we assume $u$ to be \textit{admissible}, that is:

\begin{equation}\label{2-1-condoncontrol}
    E\int_t^{t_1} \abs{u(s)}^m\,ds <\infty\,\forall m\in\mathbb{N}.
\end{equation}

It is implied by $U$ being compact. Under these hypotheses, equation \ref{2-1-controlledSDE} has a unique (indistinguishable) solution. 
Where does optimality play its role? We define running and terminal costs $L,\Psi$, both continuous and satisfying:

\begin{align}
    \abs{L(s,x,v)} \leq C(1+\abs{x}^k+\abs{v}^k) \\
    \abs{\Psi(s,x)} \leq C(1+\abs{x}^k)
\end{align}

for suitable $C,k>0$. We also define $\tau$ to be the exit time of $(s,x(s))$ from $Q$. We define:

\begin{equation}
    J(t,x;u) = E_{tx}\left\{\int_t^{\tau} L(s,x(s),u(s))\,ds + \Psi(\tau,x(\tau))\right\}
\end{equation}

for every initial condition $(t,x)\in Q$ and control $u$. We aim to minimize this criterion, that is:

\[\inf_{u\in\mathcal{U}}J(t,x;u).\]

This formulation is not mathematically formal enough, let us restate it. We begin by defining an infimum criterion with respect to a probability space, or more 
formally a \textit{probability system}, and then we'll take the infimum over all probability systems. 

\begin{definition}
    A reference probability system is a tuple $(\Omega,\{\mathcal{F}_s\},P,\omega)$ such that:
    
    \begin{enumerate}[label=\alph*)]
        \item $\nu=(\Omega,\mathcal{F}_{t_1},P)$ is a probability space
        \item $\{\mathcal{F}_s\}$ is a filtration on $\Omega$
        \item $w$ is an $\mathcal{F}$-adapted Brownian motion on $[t,t_1]$.
    \end{enumerate}

    We denote with $\mathcal{A}_{t\nu}$ the collection of all $\mathcal{F}$ progressively measurable (that is $\mathcal{B}([t,s])\times\mathcal{F}_s$-adapted), 
    $U$ valued processes $u$ such that condition \ref{2-1-condoncontrol} holds on $[t,t_1]$.
\end{definition}

We define:

\begin{equation}\label{2-1-systemopt}
    V_{\nu} = \inf_{u\in\mathcal{A}_{t\nu}} J(t,x;u)
\end{equation}

while we define:

\begin{equation}\label{2-1-opt}
    V_{PM} = \inf_{\nu} V_{\nu}.
\end{equation}

Equation \ref{2-1-systemopt} and \label{2-1-opt} respectively define $\nu$-\textit{optimality} and \textit{optimality} for those control that 
satisfy them. We adapt the definition of operator $A$ to this situation by defining for every element of the control space $v$ the functional:

\begin{equation}\label{2-1-defAwithv}
    A^{v}\Phi = \Phi_t + \sum_{i=1}^nf_i(t,x,v)\Phi_{x_i} + \frac{1}{2}\sum_{i,j=1}^n a_{ij}(t,x,v)\Phi_{x_ix_j}, \Phi\in C^{1,2}_p(\overline{Q}_0)
\end{equation}

where $a=\sigma\sigma'$. As we did in the determinist case, we provide a heuristic derivation of the Hamilton-Jacobi-Bellman equation (the verification theorem), 
and then we'll formally prove it. Let us suppose that $O=\R^n$, then $J$ is:

\begin{equation}
    J(t,x;u) = \int_t^{t_1} L(s,x(s),u(s))\,ds + \Phi(t_1,x(t_1)).
\end{equation}

By the dynamic programming principle for every $h<t_1-t$:

\[V(t,x) = \inf_{u\in\mathcal{A}}E_{tx}\left\{\int_t^{t+h} L(s,x(s),u(s))\,ds + V(t+h,x(t+h))\right\}.\]

If we take the constant control $u\equiv v$ then by Dynkin's formula we get:

\begin{align}
    0 & \leq E_{tx}V(t+h,x(t+h))-V(t,x) + E_{tx}\int_t^{t+h} L(s,x(s),v)\,ds \\
    & = E_{tx}\int_t^{t+h} A^vV(s,x(s))\,ds + E_{tx}\int_t^{t+h} L(s,x(s),v)\,ds
\end{align}

dividing by $h$ and taking the limit for $h\to0^+$:

\[0 \leq A^vV(t,x) + L(t,x,v).\]

If we take $u^{\ast}$ to be optimal, then equality holds:

\[A^{u^{\ast}}V(t,x) + L(t,x,u^{\ast}(t)) = 0.\]

We now present the verification theorem rigorously. Let us define the Hamiltonian for this problem. For every $(t,x)\in\overline{Q}_0$, 
$p\in\R^n$ and $A\in\mathcal{S}_+^n$ (set of symmetric, non-negative definite $n\times n$ matrices) we define:

\begin{equation}\label{2-1-defhamilt}
    \mathcal{H}(t,x,p,A) =  \sup_{v\in U}\left\{-f(t,x,v)\cdot p - \frac{1}{2}tr\left[a(t,x,v)\cdot A\right] - L(t,x,v)\right\}
\end{equation}

where for matrices $A,B\in\R^{n\times n}$:

\begin{equation}
    tr\left(AB\right) = \sum_{i,j=1}^n A_{ij}B_{ji}, 
\end{equation}

which is equal to $\sum_{i,j=1}^n A_{ij}B_{ij}$ for symmetric matrices.

We can now state the verification theorem using the Hamiltonian defined in \ref{2-1-defhamilt}.

\begin{theorem}
    Let $W\in C^{1,2}(Q)\cap C_p(\overline{Q})$ such that:

    \begin{align}\label{2-1-hamilcond}
        -\frac{\partial W}{\partial t}(t,x) + \mathcal{H}(t,x,D_xW,D_x^2W) = 0, & \forall(t,x)\in Q \\
        W(t,x) = \Phi(t,x), & \forall(t,x)\in\partial Q.
    \end{align}

    Then:

    \begin{enumerate}
        \item for any system $\nu$, initial condition $(t,x)\in Q$ and any $u\in \mathcal{A}_{t\nu}$ then:
        
        \begin{equation}
            W(t,x) \leq J(t,x;u)
        \end{equation}

        \item If there exists $\nu^{\ast}=(\Omega^{\ast},\{\mathcal{F}_s^{\ast}\},P^{\ast},w^{\ast})$ and $u^{\ast}\in\mathcal{A}_{t\nu^{\ast}}$ such that:
        
        \begin{equation}
            u^{\ast}(s) \in \argmin_{v\in U}\left\{f(s,x^{\ast}(s),v)\cdot D_xW(s,x^{\ast}(s)) + \frac{1}{2}tr\left[a(s,x^{\ast}(s),v)\cdot D_x^2(s,x^{\ast}(s))\right] + L(s,x^{\ast}(s),v)\right\}
        \end{equation}

        for almost all $(s,\omega)\in[t,\tau^{\ast}]\times\Omega^{\ast}$, then:

        \begin{equation}
            V_{PM}(t,x) = J(t,x;u^{\ast}).
        \end{equation}
    \end{enumerate}

    \begin{proof}
        We assume $O$ to bounded and $W\in C^{1,2}(\overline{Q})$. Because of \ref{2-1-hamilcond} for all $s\in[t,\tau]$:

        \begin{equation}\label{2-1-proofver-keyeq}
            0\leq A^{u(s)}W(s,x(s)) + L(s,x(s),u(s)).
        \end{equation}

        Because of Ito:

        \begin{equation}
            W(\tau,x(\tau)) - W(t,x) = \int_t^{\tau} A^{u(s)}W(s,x(s))\,ds + \int_t^{\tau} D_x\Phi(s,x(s))\cdot \sigma(s,x(s),u(s))\,dw(s).
        \end{equation}

        Because of estimates on SDE solution the last stochastic integral is a $\mathcal{F}_s$-martingale. Then if we take the expectation $E_{tx}$ we get:

        \begin{align}
            0 & \leq E_{tx} \int_t^{\tau}A^{u(s)}W(s,x(s))\,ds + E_{tx}\int_t^{\tau} L(s,x(s),u(s))\,ds \\
            & = E_{tx}\left(W(\tau,x(\tau)) - W(t,x)\right) - E_{tx}\int_t^{\tau} D_x\Phi(s,x(s))\cdot \sigma(s,x(s),u(s))\,dw(s) \\
            & + E_{tx}\int_t^{\tau} L(s,x(s),u(s))\,ds \\
            & = -W(t,x) + E_{tx}\left\{\int_t^{\tau} L(s,x(s),u(s))\,ds + W(\tau,x(\tau))\right\} \\
            & = -W(t,x) + J(t,x;u).
        \end{align}

        If $O$ is unbounded we define for every $\rho>0$ such that $\rho^{-1}<t_1-t_0$ the set:

        \begin{equation}
            O_{\rho} = O\cap\left\{\abs{x}<\rho\,|\,d(x,\partial O)>\frac{1}{\rho}\right\},\, Q_{\rho} = [t_0,t_1-\rho^{-1}]\times O_{\rho}
        \end{equation}

        and $\tau_{\rho}$ the exit time from $Q_{\rho}$. Then $Q_{\rho}$ is bounded, and $W\in C^{1,2}(\overline{Q}_{\rho})$, then:

        \begin{equation}
            W(t,x) \leq E_{tx}\left\{\int_t^{\tau_{\rho}} L(s,x(s),u(s))\,ds + W(\tau_{\rho},x(\tau_{\rho}))\right\}.
        \end{equation}

        We now take $\rho\to+\infty$ and get the thesis. We have convergence in probability for $\tau_{\rho}\xrightarrow{\rho\to+\infty}\tau$. 
        We prove uniform integrability of the rhs and therefore get $L^1$ convergence. We have:
        
        \begin{align}
            E_{tx}\int_t^{\tau_{\rho}} \abs{L(s,x(s),u(s))}\,ds & \leq E_{tx}\int_t^{t_1}\abs{L(s,x(s),u(s))}\,ds \\
            & \leq E_{tx}\int_t^{t_1}\left(1+\abs{x(s)}^k+\abs{u(s)}^k\right)\,ds < +\infty    
        \end{align}
        
        because $u$ is admissible and estimates on SDE solutions. While we have:

        \begin{align}
            E_{tx}\abs{W(\tau_{\rho},x(\tau_{\rho}))}^{\alpha} & \leq K E_{tx}\left(1+\abs{x(\tau_{\rho})}^k\right)^{\alpha} \\
            & \leq 2^{\alpha-1}K\left(1+E_{tx}\norm{x}^{\alpha k}\right) \leq C
        \end{align}

        for $\alpha>\frac{1}{k}$ and estimates on SDE solutions. Therefore we get:

        \begin{equation}
            \lim_{\rho\to+\infty} E_{tx}\left\{\int_t^{\tau_{\rho}} L(s,x(s),u(s))\,ds + W(\tau_{\rho},x(\tau_{\rho}))\right\} = E_{tx}\left\{\int_t^{\tau} L(s,x(s),u(s))\,ds + W(\tau,x(\tau))\right\}.
        \end{equation}

        Part $b)$ comes from equality in equation \ref{2-1-proofver-keyeq}.
    \end{proof}
\end{theorem}

% THEN DYNAMIC PROGRAMMING APPROACH AND VERIFICATION THEOREM: i present the formal derivation via $A$. Then the theorem. 
% THEN CONTROLLED MARKOV PROCESSES.

% I HAVE TO SEARCH FOR A MORE GENERAL VERSION OF EXISTENCE THEOREM for markov control policies.


% THEN EXAMPLE IF TIME ALLOWS.

% THEN REREAD, REPEAT. I DONT THINK I'LL DO SLIDES. MAYBE THEY CAN BE A GOOD ASSET. LESS INFO, LESS THINGS TO KNOW PERFECTLY. BUT I HAVE TO KNOW EVERYTHING!!!

\section{Stochastic Maximum Principle}

The stochastic analogous of Pontryagin's principle is the stochastic maximum principle. It relies on the notion of 
backward stochastic differential equation, whose solution will provide a necessary condition on the controlled system. 

\subsection{Backward Stochastic Differential Equation}

A backward stochastic differential equation is a SDE where the initial date is replaced by a final distribution. We 
start by defining a formal concept of solution and provide a general result about existence and uniqueness. 

We work with a filtered probability space $(\Omega,\mathcal{F}, P, \left\{\mathcal{F}_s\right\})$ and a Brownian motion $(w_s)$ adapted 
to the space filtration. We assume that the filtration is the natural one associated to $w$. A BSDE has the form:

\begin{equation}\label{2-3-defBSDE}
    \begin{cases}
        - dy_s = f(s,y_s,z_s)ds - z_sdw_s & \forall s\in[t,t_1]\\
        y_{t_1} = \xi, 
    \end{cases}
\end{equation}

where $f$ is real valued and $\xi$ a suitable random variable. Further conditions on $f$ and $\xi$ will be imposed 
by the existence theorem. The above definitions has to be intended in integral form. To do so we have to specify some integrability conditions. Let us define the space:

\begin{equation}\label{2-3-defS}
    \mathbb{S}^2(t,t_1) = \left\{(X_s)_{t\in[t,t_1]}\,|\, X_s\in\R \text{ is progressively measurable},\, E\left[\sup_{s\in[t,t_1]}\abs{X_s}^2\right]<+\infty\right\}
\end{equation}

and:

\begin{equation}\label{2-3-defH}
    \mathbb{H}^2(t,t_1)^n = \left\{(X_s)_{t\in[t,t_1]}\,|\, X_s\in\R^n \text{ is progressively measurable},\, E\left[\int_t^{t_1} \abs{Y_s}^2\,ds\right]<+\infty\right\}.    
\end{equation}

We can now define the solution concept of \ref{2-3-defBSDE}.

\begin{definition}
    A solution of \ref{2-3-solBSDE} is a couple $(y,z)\in\mathbb{S}^2(t,t_1)\times\mathbb{H}^2(t,t_1)^n$ such that:
    
    \[y_s = \xi + \int_s^{t_1} f(r,y_r,z_r)\,dr - \int_s^{t_1} z_r\,dw_r\] 

    holds for all $s\in[t,t_1]$. 
\end{definition}

Some measurability and integrability conditions are necessary for equation \ref{2-3-defBSDE} to make sense. Existence 
of a solution will be obtained through a classical fixed point method, which will rise from Lipschitz condition on $f$. 
We denote by $m$ the Lebesgue measure on $[t,t_1]$.

\begin{theorem}
    Let $f:\Omega\times[t,t_1]\times\R\times\R^n\rightarrow \R$ such that $f(\cdot,\cdot,y,z)$ is progressively measurable for all 
    $(y,z)\in\R^{n+1}$, $f(\cdot,\cdot,0,0)\in\mathbb{H}^2(t,t_1)^1$ and there exists $C>0$ such that:

    \begin{equation}
        \abs{f(s,y_1,z_1)-f(s,y_2,z_2)} \leq C(\abs{y_1-y_2} + \abs{z_1-z_2}),\,\forall y_1,y_2,z_1,z_2,\, m\otimes P\,a.s.
    \end{equation}

    Then for every $\xi\in L^2$ the BSDE \ref{2-3-defBSDE} has a unique solution.

    \begin{proof}
        We use completeness of the space $(\mathbb{S}(t,t_1)\times\mathbb{H}^2(t,t_1)^n, \norm{\cdot}_{\beta})$ where:

        \begin{equation}
            \norm{(y,z)}_{\beta} = \left(E\left[\int_t^{t_1} e^{\beta s}\left(\abs{y_s}^2+\abs{z_s}^2\right)ds\right]\right).
        \end{equation}

        This result is proved in the appendix. We call the previous Banach space $(X,\norm{\cdot})$. We construct the map 
        $\Phi:X\rightarrow X$ defined as $\Phi(u,v)=(y,z)$. The processes $y$ and $z$ are defined as follows. We define:

        \begin{equation}
            M_s = E\left[\xi + \int_0^{t_1} f(r,u_r,v_r)\,dr\,|\,\mathcal{F}_s\right].
        \end{equation}

        It is a square integrable martingale because:
        
        \begin{align}
            E\left\{E^2\left[\xi + \int_t^{t_1} f(r,u_r,v_r)\,dr\,|\,\mathcal{F}_s\right]\right\} & 
            % \leq C E^2\left\{E\left[\xi + \int_t^{t_1} f(r,u_r,v_r)\,dr\,|\,\mathcal{F_s}\right]\right\} \\
            % & = C E^2\left[\xi + \int_t^{t_1} f(r,u_r,v_r)\,dr\right] \\
             \leq C_0 E^2\xi + C_1E^2\int_t^{t_1} f(r,0,0)\,dr \\
            & + C_2E^2\int_t^{t_1} f(r,u_r,v_r)-f(r,0,0)\,dr \\
            & + C_3E\xi^2E\int_t^{t_1}f(r,0,0)^2\,dr \\
            & + C_4E\xi^2E\int_t^{t_1}(f(r,u_r,v_r)-f(r,0,0))^2\,dr, 
        \end{align}

        which is finite because of the assumptions and given $(u,v)\in X$. It is a martingale. Therefore, by 
        Ito's martingale representation theorem there exists a unique $z\in\mathbb{H}^2(t,t_1)^n$\footnote{Unique with respect to $\norm{\cdot}_{\mathbb{H}}$.} and $m_t\in L2$ such that:
        
        \begin{equation}
            m_s = m_t + \int_t^{s} z_r\,dw_r.
        \end{equation}

        We define $(y,z)$ as $z$ being the unique process of the martingale representation of $m_s$ while $y$ to be:

        \begin{equation}
            y_s = E\left[\xi + \int_s^{t_1} f(r,u_r,v_r)\,dr\,|\,\mathcal{F}_s\right] = m_s - \int_t^s f(r,u_r,v_r)\,dr.
        \end{equation}

        We know that $z\in\mathbb{H}^2$. We show $y\in\mathbb{S}^2$:

        \begin{align}
            E\left[\sup_{s\in[t,t_1]}\abs{y_s}^2\right] & \leq C_0E\left[\sup_{s\in[t,t_1]}\abs{m_s}^2\right] + C_1E\left[\sup_{s\in[t,t_1]}\abs{\int_t^s f(r,u_r,v_r)\,dr}^2\right],
        \end{align}

        but as before the second addend converges while for the first one:

        \[E\left[\sup_{s\in[t,t_1]}\abs{m_s}^2\right] \leq C_0 Em_t^2 + C_1E^{1/2}m_t^2E^{1/2}\sup_{s\in[t,t_1]}\left[\int_t^{s}z_r\,dw_r\right]^2+C_3E^{1/2}\sup_{s\in[t,t_1]}\left[\int_t^{s}z_r\,dw_r\right]^2,\]

        which converges beacuse Doob's inequality implies:

        \[E\left[\sup_{s\in[t,t_1]}\int_t^sz_r\,dw_r\right]^2 \leq 4E\left[\int_t^{t_1}z_r^2\,dw_r\right]<+\infty.\]

        If we prove $\Phi$ to be a contraction we'll have a unique (in $(X,\norm{\cdot})$) fixed point, therefore the thesis. 
        
        Let $(U_1,V_1),(U_2,V_2)\in X$, $(X_1,Y_1),(X_2,Y_2)$ their images and $\overline{U},\overline{V},\overline{X},\overline{Y},\overline{f}_t$ the differences between subscript $1$ and $2$. 
        By Ito's formula we have:

        \begin{align}\label{2-3-derivat}
            \abs{\overline{y_t}^2} & = -\int_t^{t_1} \frac{d}{dr}\left(e^{\beta r}\overline{y_r}^2\right) \,dr = - \int_t^{t_1} \beta e^{\beta r}\overline{y_r}^2 + e^{\beta r}\frac{d}{dr}\overline{y_r}^2 \,dr\\
            & = - \int_t^{t_1} \beta e^{\beta r}\overline{y_r}^2 \,dr - \int_t^{t_1} e^{\beta r}\left(2\overline{y_r}\right)\,d\overline{y_r} - \int_t^{t_1} e^{\beta r}\,d[\overline{y},\overline{y}]_r \\
            & = - \int_t^{t_1} \beta e^{\beta r}\overline{y_r}^2 \,dr + 2\int_t^{t_1} e^{\beta r}\overline{y_r}\overline{f}_r\,dr - 2\int_t^{t_1} e^{\beta r}\overline{y_r}\,dw_r - \int_t^{t_1} e^{\beta r}\overline{z}_r^2\,dr \\
            & = - \int_t^{t_1} e^{\beta r}\left(\beta\overline{y}_r^2-2\overline{y}_r\overline{f}_r\right)\,dr - \int_t^{t_1}e^{\beta r}\overline{z}_r^2\,dr - 2\int_t^{t_1}e^{\beta r}\overline{y}_r\overline{z}_r\,dw_r.
        \end{align}

        We show that $e^{\beta r}\overline{y}_r\overline{z}_r\in\mathbb{H}^2(t,t_1)^1$, which will guarantee the integral to vanish under expectation. We have:

        \color{red}{\[E\left[\left(\int_t^{t_1} e^{2\beta r}\overline{y}_r^2\overline{z}_r^2\,dr\right)^{1/2}\right]\leq \frac{e^{\beta t_1}}{2}E\left[\sup_{s\in[t,t_1]}\overline{y}_s^2 + \int_t^{t_1}\overline{z}_r^2\,dr\right] <+\infty.\]}
        
        \color{black}{}
        By taking the expectation on \ref{2-3-derivat}:
        
        \begin{align}
            E\overline{y}_t^2 + E\left[\int_t^{t_1}e^{\beta r} \left(\beta\overline{y}_r^2+\overline{z}_r^2\right)\,dr\right] & = 2E\left[\int_t^{t_1}e^{\beta r}\overline{y}_r\overline{f}_r\,dr\right] \\
            & \leq 2CE\left[\int_t^{t_1}e^{\beta r}\overline{y}_r\left(\overline{u}_r+\overline{v}_r\right)\,dr\right] \\
            & \leq 2CE\left[\int_t^{t_1}e^{\beta r}\overline{y}_r^2\,dr\right] + CE\left[\int_t^{t_1}e^{\beta r}\left(\overline{u}_r^2+\overline{v}_r^2\right)\,dr\right], \\
        \end{align}

        which implies:

        \[E\left[\int_t^{t_1}e^{\beta r} \left(\overline{y}_r^2+\overline{z}_r^2\right)\,dr\right] \leq E\left[\int_t^{t_1}e^{\beta r} \left(\overline{u}_r^2+\overline{v}_r^2\right)\,dr\right]\]
    \end{proof}
\end{theorem}

The previous proof used the concept of local martingale and the Burkholder-Davis-Gundy inequality, they are detailed in the appendix. In particular the fact that if $X\in\mathbb{H}^2$ then $\int X\,dw$ is a martingale uses these concepts.

A notable case is the one with a generator $f$ linear in $y$ and $z$. That is, there exists 
$a,b$ bounded progressively measurable processes valued in $\R$ and $\R^n$ and $c\in\mathbb{H}(t,t_1)^1$ which define:

\begin{equation}\label{2-3-defBSDE-linear}-dy_s = \left(a_sy_s + z_sb_s + c_s\right)ds - z_sdw_s,\,y_{t_1}=\xi.\end{equation}

\begin{proposition}
    The unique solution $(y,z)$ of \ref{2-3-defBSDE-linear} is:
    
    \begin{equation}\label{2-3-linearfsol}
        \Gamma_sy_s = E\left[\Gamma_{t_1}\xi + \int_s^{t_1} \Gamma_rc_r\,dr\,|\,\mathcal{F}_s\right],
    \end{equation}

    and $z_s$ is defined via the Martingale representation of \ref{2-3-linearfsol}. The process $\Gamma$ is defined by:

    \[d\Gamma_s = \Gamma_s\left(a_sds + b_sdw_s\right),\,\Gamma_t=1.\]
\end{proposition}

\subsection{Stochastic Maximum Principle}

The stochastic counterpart of Pontryagin's principle is the stochastic maximum principle. We consider the controlled diffusion process $x$ on $\R^n$ defined by:

\begin{equation}
    dx_s = f(r,x(r),u(r))dr + \sigma(r,x(r),u(r))dw(r)
\end{equation}

and let the functional to be maximized:

\begin{equation}\label{2-3-newfunctmax}
    J(t,x;u) = E_{tx}\left\{\int_t^{t_1}L(s,x(s),u(s))\,ds + \Phi(t_1,x(t_1))\right\},
\end{equation}

where the running and terminal cost satisfy the usual assumptions. We aim at maximizing functional $J$ over all 
admissible systems, that is:

\[\inf_{\nu}\inf_{u\in\mathcal{A}_{t\nu}} J(t,x;u).\]

We consider again the value function related to this problem, that is a function $V$, with suitable smoothness conditions, such that:

\begin{equation}
    -\frac{\partial V}{\partial t} + \mathcal{H}(t,x,D_xV,D_x^2V) = 0,
\end{equation}

where:

\begin{equation}
    \mathcal{H}(s,x,v,p,A) = - f(s,x,v)\cdot p - \frac{1}{2}tr\left[\sigma\sigma'(s,x,v)\cdot A\right] - L(s,x,v).
\end{equation}

As in the determinist case, under the assumption of optimality the value function $V$ will define the solution 
of a differential equation, more precisely a BSDE. The backward stochastic differential equation will be an 
analogous of the deterministic adjoint equation:

\[\begin{cases}
    \dot{p}(s) = D_xH(s,x(s),u(s),p(s)) \\
    p(t_1) = D\psi(x(t_1)),
\end{cases}\]

as stated in \ref{1-3-pontry}. In particular, we will need the functional:

\begin{equation}\label{3-2-defG}\mathcal{G}(t,x,v,y,z) = -f(s,x,v)\cdot y - tr\left[\sigma'(s,x,v)\cdot z\right] - L(s,x,v).\end{equation}

This functional will define the BSDE in the next theorem.

\begin{theorem}\label{2-3-stochmaxprinc}
    Let $u^{\ast}$ be an optimal control and $x^{\ast}$ the corresponding diffusion process, and the value function 
    $V\in C^{1,2}(O)\cap C(\overline{O})$. Then $V$ satisfies:
    
    \begin{equation}\label{2-3-maxcond}
        \mathcal{H}(s,x_s^{\ast}, u_s^{\ast}, D_xV(s,x^{\ast}_s), D_x^2V(s,x^{\ast}_s)) = \sup_{v \in U}\mathcal{H}(s,x_s^{\ast}, v, D_xV(s,x^{\ast}_s), D_x^2V(s,x^{\ast}_s)),
    \end{equation}

    and the pair $(y_s,z_s)=\left(D_xV(s,x_s^{\ast}),D_x^2v(s,x_s^{\ast})\sigma(s,x_s^{\ast},u_s^{\ast})\right)$ solves the BSDE:

    \begin{equation}\label{2-3-maxprinc-BSDE}
        -dy_s = D_x\mathcal{G}(s,x_s^{\ast}, u_s^{\ast}, y_s, z_s)ds - z_sdw_s,
    \end{equation}

    with final condition:

    \begin{equation}\label{2-3-maxprinc-finalcond}
        y_{t_1} = D_x\Psi(t_1,x_{t_1}).
    \end{equation}

    \begin{proof}
        We drop the $^{\ast}$. We consider the optimal system over which the control is defined. 
        Since $u$ is optimal we have:

        \begin{align}
            V(s,x(s)) & = E_{sx(s)}\left[\int_s^{t_1} L(r,x(r),u(r))\,dr + \Psi(t_1,x(t_1))\right] \\ 
            & = E_{tx}\left[\int_t^{t_1} L(r,x(r),u(r))\,dr + \Psi(t_1,x(t_1))\right] - \int_t^s L(r,x(r),u(r))\,dr.  
        \end{align}

        We then apply Ito's formula and get:

        \begin{align}
            \int_t^{t_1} \partial_sV(r,x(r)) + D_xV(r,x(r))f(r,x(r),u(r)) + \frac{1}{2}tr\left[\sigma\sigma'(r,x(r),u(r))D_xx^2V(r,x(r))\right] \,dr \\
            + \int_t^{t_1} D_xV(r,x(r))\sigma(r,x(r),u(r))\,dw_r + V(t,x) \\
            = \int_t^{t_1} -f(r,x(r),u(r))\,dr + \int_t^{t_1}\alpha_r\,dw_r + V(t,x)
        \end{align}

        where we used the integral representation of a martingale, which implies:

        \begin{equation}\label{2-3-hjbinproof}
            -\frac{\partial V}{\partial t} + \mathcal{H}(s,x(s),u(s),D_xV(s,x(s)),D_x^2V(s,x(s))) = 0.
        \end{equation}

        We then prove the BSDE using the continuity of triple spacial derivatives of the value functions and 
        the fact that \ref{2-3-hjbinproof} at $s$ has maximum in $x(s)$, which implies:

        \[\frac{\partial}{\partial x}\left(\frac{\partial V}{\partial t}(s,x) - \mathcal{H}(s,x,u(s),D_xV(s,x),D_x^2V(s,x))\right)|_{x=x(s)}=0.\]

        We compute the derivative and get:

        \begin{align*}
            \frac{\partial^2V}{\partial x\partial t}(s,x(s)) + D_xf(s,x,u(s))|_{x=x(s)}D_xV(s,x(s)) + f(s,x(s),u(s))D_x^2V(s,x(s)) \\
            + \frac{1}{2}tr\left[D_x(\sigma\sigma')(s,x,u(s))|_{x=x(s)}\cdot D_x^2V(s,x(s))\right] + \frac{1}{2}tr\left[\sigma\sigma'(s,x(s),u(s))\cdot D_x^3V(s,x(s))\right] + D_xL(s,x(s),u(s)) \\
            = \frac{\partial^2V}{\partial x\partial t}(s,x(s)) + f(s,x(s),u(s))D_x^2V(s,x(s)) + \frac{1}{2}tr\left[\sigma\sigma'(s,x(s),u(s))\cdot D_x^3V(s,x(s))\right] \\
            + f(s,x(s),u(s))D_xV(s,x(s)) + tr\left[D_x\sigma'(s,x(s),u(s))\cdot D_x^2V(s,x(s))\sigma(s,x(s),u(s))\right] + D_xL(s,x(s),u(s)),
        \end{align*}

        which implies:

        \begin{equation}\label{2-3-relD3G}
            \begin{align*}\frac{\partial^2V}{\partial x\partial t}(s,x(s)) + f(s,x(s),u(s))D_x^2V(s,x(s)) + \frac{1}{2}tr\left[\sigma\sigma'(s,x(s),u(s))\cdot D_x^3V(s,x(s))\right] \\
            = -D_x\left[\mathcal{G}(s,x(s),u(s),D_xV(s,x(s)),D_x^2V(s,x(s))\sigma(s,x(s),u(s)))\right].
            \end{align*}
        \end{equation}

        Equation \ref{2-3-relD3G} implies that $y_s$ is the unique solution to the BSDE \ref{2-3-maxprinc-BSDE}:

        \begin{align*}
            -dy_s & = -d[D_xv(s,x(s))] =  -\frac{\partial^2V}{\partial x\partial t}(s,x(s))ds - f(s,x(s),u(s))D_x^2V(s,x(s))ds \\
            & - D_x^2V(s,x(s))\sigma(s,x(s),u(s))dw_s + \frac{1}{2}tr\left[\sigma\sigma'(s,x(s),u(s))\cdot D_x^3V(s,x(s))\right]ds \\
            & = D_x\left[\mathcal{G}(s,x(s),u(s),D_xV(s,x(s)),D_x^2V(s,x(s))\sigma(s,x(s),u(s)))\right]ds - D_x^2V(s,x(s))\sigma(s,x(s),u(s))dw_s,
        \end{align*}

        and:

        \[V(t_1,x(t_1)) = E_{t_1x(t_1)}\left[\int_{t_1}^{t_1} L(r,x(r),u(r))\,dr + \Psi(t_1,x(t_1))\right] = \Psi(t_1,x(t_1)), \]

        which implies:

        \[D_xV(t_1,x(t_1)) = D_x\Psi(t_1,x(t_1)).\]
    \end{proof}
\end{theorem}

\newpage

\chapter{Viscosity Solutions}

We now study viscosity solutions and their relation with optimal control.

\section{Introduction}

As previously mentioned, in many instances the value function arising from an optimal control problem may fail to be continuously 
differentiable. If that happens the derivation of the Hamilton-Jacobi equation is no longer valid, but more importantly 
the notion of classical solution to it does hold anymore. Therefore, we have to weaken the notion of solution in order to get 
a consistent and unique solution to the dynamic programming equation for non-differentiable value functions. The \textit{viscosity solution} 
is exactly what we are searching for. It arises from a standard procedure called vanishing viscosity, which allows us to compute the 
solution of a fully non-linear first order PDE as the limiting solution of quasilinear parabolic PDEs, obtained via infinitesimal 
perturbations of second order derivatives. 

\subsection{Non-differentiable value functions}

Let us consider the calculus of variation problem:

\begin{equation}
    \inf_{x\in Lip([0,1];[-1,1])} \int_t^{t_1} 1 + \frac{1}{4}(\dot{x}(s))^2 \,ds,
\end{equation}

where $Lip(I;U)$ is the collection of Lipschitz continuous functions from $I$ to $U$. The Hamiltonian related to 
this problem is:

\[H(t,x,p) = \max_{v\in [-1,1]} \left\{-v\cdot p - 1 - \frac{1}{4}v^2\right\}.\]

We can explicitly compute the Hamiltonian and get:

\[H(t,x,p)=p^2-1.\]

Then the Hamilton-Jacobi equations read:

\[\begin{cases}
    \dot{x}^{\ast}(s) = -H_p(s,x^{\ast}(s),p^{\ast}(s)) =  2p^{\ast}(s)\\
    \dot{p}^{\ast}(s) = H_x(s,x^{\ast}(s),p^{\ast}(s)) = 0,
\end{cases}\]

therefore, we get:

\[\dot{x}(s)^{\ast} = 2p^{\ast},\,s\in[0,1],\]

for some $p^{\ast}\in\R$. We now compute the exit time of $(s,x(s))=(s,2(s-t)p^{\ast}+x)$ with initial data $(t,x)$. If $p=0$ then:

\[\tau=1,\,\abs{x}<1.\]

If $p>0$ then $x(s)=2(s-t)p+x$ is increasing, which implies that the system is going to exit from the right boundary, that is from $x(s)=1$, and if 
that happens before time $s=1$ the exit time will be determined by:

\[2(s-t)p+x=1\Rightarrow s=t+\frac{1-x}{2p}.\]

$x(s)=1$ for $s<1$ if:

\[2(1-t)p+x\geq1\Rightarrow p\geq t+\frac{1-x}{2p},\]

therefore:

\[\tau=\begin{cases}
    1 & p\geq t+\frac{1-x}{2p} \\
    t+\frac{1-x}{2p} & p> t+\frac{1-x}{2p}.
\end{cases}\]

Analogously, if $p<0$:

\[\tau=\begin{cases}
    1 & p\leq t-\frac{1+x}{2p} \\
    t-\frac{1+x}{2p} & p< t-\frac{1+x}{2p}.
\end{cases}.\]

We now solve:

\[\inf_{p\in\R} \int_t^{\tau}1+p^2\,ds=\begin{cases}
    (1+p^2)(1-t), & p=0\text{ or }p>0\land p\geq \frac{1-x}{2(1-t)}\text{ or }p<0\land p\leq\frac{-1-x}{2(1-t)}\\
    (1+p^2)\frac{1-x}{2p}, & p>0\land p\geq \frac{1-x}{2(1-t)} \\
    -(1+p^2)\frac{1+x}{2p}, & p\leq\frac{-1-x}{2(1-t)},
\end{cases}\]

which is solved as follows:

\[V(t,x)=\begin{cases}
    1-t & \abs{x}\leq t \\
    1-t& \abs{x}\geq t,
\end{cases}\]

which is continuous on the whole space, but clearly no differentiable in $\abs{x}=t$.

\subsection{Vanishing viscosity}

We now euristichally expose a technique called vanishing viscosity, which is widely used in 
calculus of variations problems and will show us one of the origins of the viscosity 
solution notion. Let us consider the initial value problem:

\begin{equation}\label{4-1-fullynonlinear}
    \begin{cases}
        u_t + H(u,Du) = 0,& \R^n\times(0,+\infty) \\
        u = g,& \R^n\times\{t=0\}.
    \end{cases}
\end{equation}

The method of characteristics shows that there cannot be a smooth solution of the above problem over the whole positive 
real line. Indeed, a weaker notion of solution is needed. One approach is to use Hopf-Lax solution concept. We are not 
interested in it, instead we start by perturbing the system as:

\begin{equation}
    \begin{cases}
        u^{\epsilon}_t + H(u^{\epsilon},Du^{\epsilon}) -\epsilon\Delta u^{\epsilon}= 0,& \R^n\times(0,+\infty) \\
        u^{\epsilon} = g,& \R^n\times\{t=0\},
    \end{cases}
\end{equation}

so that the fully non-linear system in \ref{4-1-fullynonlinear} becomes a semilinear one, which turn out to have a smooth solution. 
Then, we take $\epsilon\to0$. We expect the solution $u^{\epsilon}$ to lose the bounds on the derivatives, as they strongly depend on the 
regularization effect of $\epsilon\Lambda$. Turns out that many times Ascoli-Arzela theorem's hypotheses are satisfied, that is $(u^{\epsilon})_{\epsilon}$ 
is uniformly bounded and equicontinuous, then we have local uniform convergence along a subsequence $u^{\epsilon_j}$. We now use the limit
$u\xleftarrow{j\to+\infty}u^{\epsilon_j}$ as a solution. 
We now it to be continuous but we lack information about its derivatives. We will then verify these information using test functions. 
Unlike the classical variational weak solution concept, where integration by part play the central role, we will use the maximum principle 
to translate the derivates of $u$ onto the test functions. 

Let us take $v\in C^{\infty}(\R^n\times(0,+\infty))$ and suppose that $u-v$ has a strict local maximum at $(x_0,t_0)$, then:

\[(u-v)(x_0,t_0) > (u-v)(x,t),\]

for all $(t,x)$ sufficiently close to $(x_0,t_0)$. It can be shown that it implies that there exists $J>0$ such that for all 
$j>J$ there exists $(x_{\epsilon_j},t_{\epsilon_j})$ such that:

\[(u^{\epsilon}-v)(x_{\epsilon_j},t_{\epsilon_j}) \geq (u^{\epsilon}-v)(x,t),\]

for $(x,t)$ sufficiently close to $(x_{\epsilon_j},t_{\epsilon_j})$ and such that:

\[(x_{\epsilon_j},t_{\epsilon_j})\xrightarrow{j\to+\infty}(x_0,t_0).\]

Because $(u^{\epsilon}-v)$ has a local maximum at $(x_{\epsilon_j},t_{\epsilon_j})$:

\[Du^{\epsilon_j}(x_{\epsilon_j},t_{\epsilon_j})=Dv(x_{\epsilon_j},t_{\epsilon_j}),\,u^{\epsilon_j}_t(x_{\epsilon_j},t_{\epsilon_j}) =v(x_{\epsilon_j},t_{\epsilon_j}),\]

and:

\[-\Lambda u^{\epsilon_j}(x_{\epsilon_j},t_{\epsilon_j}) \geq -\Lambda v(x_{\epsilon_j},t_{\epsilon_j}).\]

Therefore, we get:

\[v_t(x_{\epsilon_j},t_{\epsilon_j}) +H(Dv(x_{\epsilon_j},t_{\epsilon_j}) ,x_{\epsilon_j})\leq \Lambda v(x_{\epsilon_j},t_{\epsilon_j}) \xrightarrow{j\to+\infty}0\]

Analogous computations can be done for local minimum of $u-v$, obtaining the opposite inequality above.

We can now grasp the intuition behind the following definition.

\begin{definition}
    A viscosity solution of \ref{4-1-fullynonlinear} is a function $u$ bounded and uniformly continuous on $\R^n\times[0,T]$ for all $T>0$ such that 
    for all $v\in C^{+\infty}(\R^n\times(0,+\infty))$:

    \[v_t(x,t)+H(Dv(x,t),x)\leq0\]

    for all $(x,t)\in\arg \max\{u-v\}$ and:

    \[v_t(x,t)+H(Dv(x,t),x)\geq0\]

    for all $(x,t)\in\arg \min\{u-v\}$. Furthermore, $u\equiv g$ for $t=0$. 
\end{definition}

\subsection{Abstract dynamic programming and viscosity solutions}

We now present an abstraction of the dynamic programming principle, which will allow us to define the viscosity solutions of the dynamic programming 
equation. Let $\Sigma$ be a closed subset of a Banach space and $\mathcal{C}$ a collection of functions on $\Sigma$, closed under addition:

\[\phi,\psi\in\mathcal{C}\Rightarrow \phi+\psi\in\mathcal{C}.\]

We consider the family of operators $\{\mathcal{T}_{tr}\}_{t_0\leq t\leq r\leq t_1}$ such that: 

\begin{equation}\label{4-1-identityofT}
    \mathcal{T}_{tt}\phi = \phi,\, \forall\phi\in\mathcal{C},
\end{equation}

\begin{equation}\label{4-1-leqT}
    \mathcal{T}_{tr}\phi\leq\mathcal{T}_{ts}\psi \text{ if } \phi\leq \mathcal{T}_{rs}\psi,
\end{equation}

and:

\begin{equation}\label{4-1-geqT}
    \mathcal{T}_{tr}\phi\geq\mathcal{T}_{ts}\psi \text{ if } \phi\geq \mathcal{T}_{rs}\psi.
\end{equation}

Conditions \ref{4-1-leqT} and \ref{4-1-geqT} are a weaker version of monotonicity; they imply it together with \ref{4-1-identityofT}. 
Moreover, they also imply the semigroup property, provided that $\mathcal{T}_{rt}:\mathcal{C}\rightarrow\mathcal{C}$. Under this assumption,
the two conditions are equivalent to monotonicity. The semigroup property:

\begin{equation}
    \mathcal{T}_{tr}\left(\mathcal{T}_{rs}\psi\right)=\mathcal{T}_{ts}\psi,\,\mathcal{T}_{rs}\psi\in\mathcal{C},
\end{equation}

is going to be the dynamic programming principle. Indeed, let us consider the classical optimal control problem defined on a bounded set 
$O\subset\R^n$, which we set to be $\Sigma=\overline{O}$ and $\mathcal{C}=\mathcal{M}(\Sigma)$, the collection of measurable functions bounded by below. Then as in chapter 1 we aim at minimize a functional, 
we set this functional to be the operator $\mathcal{T}$. Let us define:

\begin{equation}\label{4-1-valuefunctreform}
    \mathcal{T}_{t,r;u}\psi(x) = \int_t^{\tau\wedge r}L(s,x(s),u(s)),\,ds + g(\tau,x(\tau))\chi_{\tau<r} + \psi(x(r))\chi_{\tau\geq r}, 
\end{equation}

which gives:

\begin{equation}\label{4-1-defofT_dynam}
    \mathcal{T}_{tr}\psi = \inf_{u\in\mathcal{U}(t,x)} \mathcal{T}_{t,r;u}\psi.
\end{equation}

Under the usual assumption on the running and terminal costs, as well as on the control space $U$ we now that the value function defined in ref{4-1-valuefunctreform} 
is measurable and bounded by below. Therefore, $\mathcal{T}_{rt}:\mathcal{C}\rightarrow\mathcal{C}$ and we can formulate the semigroup property just by asking $\psi\in\mathcal{C}$. 
It is clear by its definition that the dynamic programming principle is translated as:

\[\mathcal{T}_{tt_1}\psi(x)=\mathcal{T}_{tr}\left(\mathcal{T}_{rt_1}\psi\right)(x),\]

for $(t,x)\in\overline{Q}$ and $\psi\in\mathcal{C}$. We now derive the abstract dynamic programming equation. The same procedure as in chapter 1 gives:

\[-\frac{1}{h}\left[\mathcal{T}_{tt+h}V(t+h,\cdot)(x)-V(t,x)\right]=0.\]

What happens if we let $h\to0$? We ask for the existence of a family of non-linear operators which will play the role of the Hamiltonian. 
Let $\Sigma'\subset\Sigma$ and $\mathcal{D}\subset C([t_0,t_1)\times\Sigma')$ and $\{\mathcal{G}_t\}_{t\in[t_0,t_1]}$ functions on $\Sigma$ such that:

\begin{equation}\label{4-1-conditiononG_t}
    \lim_{h\to 0}\frac{1}{h}\left[\mathcal{T}_{tt+h}V(t+h,\cdot)(x)-V(t,x)\right] = \frac{\partial}{\partial t}w(t,x) - (\mathcal{G}_tw(t,\cdot))(x)
\end{equation}

for all $w\in\mathcal{D},(t,x)\in Q=[t_0,t_1)\times\Sigma'$. The space $\mathcal{D}$ is such that:

\begin{equation}\label{4-1-conditiononD}
    \forall w\in\mathcal{D} \text{ the functions }\frac{\partial w}{\partial t},\,\mathcal{G}_tw(t,\cdot) \text{ are continuous on $Q$ and }w(t,\cdot)\in\mathcal{C} \forall t\in[t_0,t_1], 
\end{equation}

and $\mathcal{D}$ is a vector space:

\begin{equation}
    w,\tilde{w}\in\mathcal{D} \Rightarrow w+\tilde{w}\in\mathcal{D},\,\lambda w\in\mathcal{D}.
\end{equation}

The elements of $\mathcal{D}$ are called test functions and $\mathcal{G}_t$ the infinitesimal generator of $\mathcal{T}_{tr}$. Explicit choices of $\mathcal{C}$ and 
$\mathcal{D}$ will vary from case to case, usually they are chosen to satisfy certain integrability conditions on the functions. 

If we require the existence of a test function space and an infinitesimal generator of the semigroup given by the value function, the 
dynamic programming equation becomes:

\begin{equation}\label{4-1-newprogreq}
    -\frac{\partial }{\partial t}V(t,x) + (\mathcal{G}_tV(t,\cdot))(x) = 0,\,(t,x)\in Q.
\end{equation}

Then if a function $V\in\mathcal{D}$ satisfies \ref{4-1-newprogreq} point wise it is called a classical solution of it. Thanks to this reformulation we will be able 
to weaken this notion of solution and finally get the viscosity solution of the dynamic programming equation.

We point out that in the canonical deterministic optimal control studied in chapter 1, the infinitesimal generator has the form:

\[(\mathcal{G}_t\phi)(x) = H(t,x,D\phi(x)) = \sup_{v \in U}\left\{-f(t,x,v)\cdot D\phi(x) - L(t,x,v)\right\},\]

with test functions space $\mathcal{D}= C^1(Q)\cap \mathcal{M}(\overline{Q})$ and $\Sigma'=O$.

In view of what we saw in previous sections, we give the definition of viscosity solution. 

\begin{definition}\label{4-1-defviscsol}
    Let $W\in C([t_0,t_1]\times\Sigma)$. Then:
    \begin{enumerate}
        \item $W$ is a \textit{viscosity subsolution} of \ref{4-1-newprogreq} in $Q$ if for every $w\in\mathcal{D}$:
        \begin{equation}\label{4-1-defsubsolution}
            -\frac{\partial}{\partial t}w\left(\overline{t},\overline{x}\right) + \left(\mathcal{G}_{\overline{t}}w\left(\overline{t},\cdot\right)\right)\left(\overline{x}\right) \leq 0,
        \end{equation} 
        at every:
        \[\left(\overline{t},\overline{x}\right)\in \arg \max_{(t,x)\in Q}\left\{(W-w)(t,x)\right\},\]

        and $W\left(\overline{t},\overline{x}\right)=w\left(\overline{t},\overline{x}\right)$.
        \item $W$ is a \textit{viscosity supersolution} of \ref{4-1-newprogreq} in $Q$ if for every $w\in\mathcal{D}$:
        \begin{equation}\label{4-1-defsupersolution}
            -\frac{\partial}{\partial t}w\left(\overline{t},\overline{x}\right) + \left(\mathcal{G}_{\overline{t}}w\left(\overline{t},\cdot\right)\right)\left(\overline{x}\right) \geq 0,
        \end{equation} 
        at every:
        \[\left(\overline{t},\overline{x}\right)\in \arg \min_{(t,x)\in Q}\left\{(W-w)(t,x)\right\},\]

        and $W\left(\overline{t},\overline{x}\right)=w\left(\overline{t},\overline{x}\right)$.
        \item $W$ is a \textit{viscosity solution} if it is both a subsolution and a supersolution.
    \end{enumerate}
\end{definition}

As a first step after the definition of viscosity solutions, we prove its consistency with the classical notion. To do so we require the 
operator to satisfy a maximum principle, that is if $\mathcal{G}_t$ is a general 
operator and:

\[\mathcal{D}=\left\{W\in C([t_0,t_1]\times\Sigma)\,|\,W_t(t,x),(\mathcal{G}_tW(t,\cdot))(x)\in C(Q)\right\},\]

then $\mathcal{G}_t$ satisfies the \textit{maximum principle} if:

\[\mathcal{G}_t\phi(\overline{x})\geq\mathcal{G}_t\psi(\overline{x})\]

for every $\overline{x}\in\arg \max \{(\phi-\psi)(x)\,|\,x\in\Sigma\}\cap\Sigma'$ with $\phi(\overline{x})=\psi(\overline{x})$. Now,
if $W$ is a classical solution of \ref{4-1-newprogreq}, then for a $w\in\mathcal{D}$ and $\left(\overline{t},\overline{x}\right) \in \arg \max_{(t,x)\in Q}\left\{(W-w)(t,x)\right\}$ 
and $W\left(\overline{t},\overline{x}\right)=w\left(\overline{t},\overline{x}\right)$ then:

\begin{align}
    -\frac{\partial}{\partial t}w\left(\overline{t},\overline{x}\right) + \left(\mathcal{G}_{\overline{t}}w\left(\overline{t},\cdot\right)\right)\left(\overline{x}\right) \leq -\frac{\partial}{\partial t}W\left(\overline{t},\overline{x}\right) + \left(\mathcal{G}_{\overline{t}}W\left(\overline{t},\cdot\right)\right)\left(\overline{x}\right) = 0,
\end{align}

since we asked continuity of time derivatives. For the supersolution recall that:

\[ \max \{(\phi-\psi)(x)\,|\,x\in \Sigma\} = \min \{-(\phi-\psi)(x)\,|\,x\in \Sigma\}.\]

If $\mathcal{G}_t$ is the infinitesimal generator of a two-parameter semigroup the connection between classical and viscosity solutions is even stronger.

\begin{proposition}
    Let $W\in\mathcal{D}$. Then $W$ is a classical solution of \ref{4-1-newprogreq} if and only if it is a viscosity solution of \ref{4-1-newprogreq} in $Q$.

    \begin{proof}
        If $W$ is a viscosity solution, since it is also a test function then \ref{4-1-subsolution} and \ref{4-1-supersolution} hold for every point $(t,x)\in Q$, 
        which implies that:
        \[-\frac{\partial}{\partial t}w(t,x) + \left(\mathcal{G}_tw(t,\cdot)\right)(x) = 0,\,\forall(t,x)\in Q.\] 

        If $W$ is a classical solution and prove the subsolution property. Let $w\in\mathcal{D}$ and $(t,x)$ as usual. Since $w\geq W$:

        \begin{align*}
            -\frac{\partial}{\partial t}w(t,x) + (\mathcal{G}_tw(t,\cdot))(x) & = - \lim_{h\to0} \frac{1}{h}\left[(\mathcal{T}_{tt+h}w(t+h,\cdot))(x)-w(t,x)\right] \\
            & \leq - \lim_{h\to0} \frac{1}{h}\left[(\mathcal{T}_{tt+h}W(t+h,\cdot))(x)-W(t,x)\right] \\
            & = -\frac{\partial}{\partial t}W(t,x) + (\mathcal{G}_tW(t,\cdot))(x) = 0.
        \end{align*}

        The supersolution property is proven similarly.
    \end{proof}
\end{proposition}

We now prove that the family of linear operators $\mathcal{T}_{tt_1}$ defined in \ref{4-1-defofT_dynam} defines viscosity solution of the dynamic programming 
equation \ref{4-1-newprogreq}.

\begin{theorem}\label{4-1-thm:Vtoviscosity}
    Let $\{\mathcal{T}_{tr}\}_{t_0\leq t\leq r\leq t_1}$ such that \ref{4-1-identityofT},\ref{4-1-geqT},\ref{4-1-leqT} and also there exists a vector space 
    $\mathcal{D}$ and another family of operator $\{\mathcal{G}_t\}_{t\in[t_0,t_1]}$ such that \ref{4-1-conditiononD} and \ref{4-1-conditiononG_t} hold. Let:
    
    \[V(t,x)=\left(\mathcal{T}_{tt_1}\psi\right)(x).\]

    If $V\in C(Q)$ then it is a viscosity solution of \ref{4-1-newprogreq}.

    \begin{proof}
        Let us prove the subsolution condition. Let $w\in\mathcal{D}$ and $(t,x)$ maximizer of $V-w$ in $\overline{Q}$ and $V(t,x)=w(t,x)$. Then $V\leq w$ on 
        $Q$, which implies by \ref{4-1-geqT} that we have:
        
        \begin{equation}\label{4-1-proof1-firsstineq}(\mathcal{T}_{tr}w(r,\cdot))(x) \geq (\mathcal{T}_{tt_1}\psi)(x) = V(t,x) = w(t,x).\end{equation}

        Now, because of \ref{4-1-conditiononD} and \ref{4-1-conditiononG_t} we can compute:
        
        \[-\frac{\partial}{\partial t}w(t,x) + (\mathcal{G}_tw(t,\cdot))(x) = - \lim_{h\to0} \frac{1}{h}\left[(\mathcal{T}_{tt+h}w(t+h,\cdot))(x)-w(t,x)\right]\leq0\]

        because of \ref{4-1-proof1-firsstineq}. The supersolution condition is proven analogously.
    \end{proof}
\end{theorem}

\subsection{The Standard Approach}

We have talked about \textit{viscosity solutions} for $\mathcal{G}_t$ that are infinitesimal generators of semigroups. Historically the notion of viscosity solution was introduced 
for partial differential equations, that is when $\mathcal{G}_t$ is a partial differential operator. In this section we will introduce this classical notion and provide some properties, 
establishing as well the relevant links to the already proven theory. 

Let $\mathcal{G}_t$ be a partial differential operator defined As

\begin{equation}\label{4-1-GPDEop}
    (\mathcal{G}_t\phi)(x) = F(t,x,D\phi(x), D^2\phi(x), \phi(x)),
\end{equation}

where $F$ is taken to be continuous. We assume that $\sigma'=O\subset\R^n$ open and $\sigma$ to be its closure. 
Furthermore, we assume that test functions $\mathcal{C}$ and $\mathcal{D}$ are such that

\begin{equation}\label{4-1-condonCandD}
    C_p(\overline{O})\cap\mathcal{M}(\overline{O}) \subset \mathcal{C},\,C_p^{\infty}(\overline{Q})\cap\mathcal{M}(\overline{Q}) \subset \mathcal{D},
\end{equation}

where $Q=[t_0,t_1)\times O$ as usual. A useful result is the following.

\begin{lemma}
    The operator $\mathcal{G}_t$ obeys to the maximum principle if and only if $F$ is elliptic.

    \begin{proof}
        Recall that $F$ is elliptic If

        \begin{equation}\label{4-1-defelliptic}
            F(t,x,p,A+B,V) \leq F(t,x,p,A,V)
        \end{equation}

        for all $(t,x,p,V)\in Q\times\R^n\times\R$ and all $A,B\in Sym(\R^{n\times n})$ with $B\geq0$.

        Let us assume that $F$ is elliptic, then for generic $\phi,\psi\in C^2(O)$ and $\overline{x}\in\arg \max\{(\phi-\psi)(x)|x\in\overline{O}\}\cap O$ 
        with $\phi(\overline{x})=\psi(\overline{x})$ we have that $D\phi(\overline{x})=D\psi(\overline{x})$ and $D^2\phi(\overleftarrow{x})\leq D^2\psi(\overleftarrow{x})$, 
        therefore

        \[F(t,\overline{x}, D\phi(\overline{x}), D^2\phi(\overline{x}),\phi(\overline{x}))\geq F(t,\overline{x}, D\psi(\overline{x}), D^2\psi(\overline{x}),\psi(\overline{x})),\]

        just by using \ref{4-1-defelliptic} with $B=D^\psi(\overline{x})-D^2\phi(\overline{x})\geq0$. If $F$ is not elliptic 
        for some $(t,\overline{x},p,A,V)$ then the functions

        \[\psi(x) = V+p\cdot(x-\overline{x}) + \frac{1}{2}(A+B)(x-\overline{x})\cdot(x-\overline{x}),\]

        \[\phi(x) = V+p\cdot(x-\overline{x}) + \frac{1}{2}A(x-\overline{x})\cdot(x-\overline{x}),\]

        for $x\in\overline{O}$ are such that $\overline{x}$ maximizes the distance between them, they coincide over it, but 
        the value of $F$ in $\overline{x}$ and $\phi,\psi$ does not satisfy the maximum principle, since \ref{4-1-defelliptic} 
        does not hold.
    \end{proof}
\end{lemma}

I this context we give the following definition.

\begin{definition}\label{4-1-defviscforW}
    Given a function $W\in C(\overline{O})$ we define the equation
    \begin{equation}\label{4-1-PDEwithW}
        -\frac{\partial}{\partial t}W(t,x) + F(t,x,D_xW(t,x),D_x^2W(t,x),W(t,x))=0
    \end{equation}

    \begin{enumerate}[label=(\alph*)]
        \item $W$ is a \textit{viscosity subsolution} of \ref{4-1-PDEwithW} in $Q$ if for each $w\in C^{\infty}(Q)$ then 
        
        \begin{equation}
            -\frac{\partial}{\partial t}w(t,x) + F(t,x,D_xw(t,x),D_x^2w(t,x),w(t,x))\leq0
        \end{equation}

        at every $(t,x)\in Q$ which locally maximizes $W-w$ yielding $0$.
        
        \item $W$ is a \textit{viscosity supersolution} of \ref{4-1-PDEwithW} in $Q$ if for each $w\in C^{\infty}(Q)$ then 
        
        \begin{equation}
            -\frac{\partial}{\partial t}w(t,x) + F(t,x,D_xw(t,x),D_x^2w(t,x),w(t,x))\geq0
        \end{equation}

        at every $(t,x)\in Q$ which locally minimizes $W-w$ yielding $0$.

        \item $W$ is a \textit{viscosity solution} if it is both a viscosity subsolution and supersolution.
    \end{enumerate}
\end{definition}

We will prove that definitions \ref{4-1-defviscsol} and \ref{4-1-defviscforW} are equivalent for partial differential operators 
of the form \ref{4-1-GPDEop}, provided that the function has sub-polynomial growth and is measurable.

An handy way of verifying both the definitions is the following.

\begin{lemma}
    If the conditions in the definitions \ref{4-1-defviscsol}, \ref{4-1-defviscforW} is satisfied only at strict extrema 
    then it is satisfied at every extrema.

    \begin{proof}
        We prove it for the minimum of definition \ref{4-1-defviscforW}, the other assertions are proved similarly. Let 
        $(t,x)$ be a non strict minimum of $W-w$ such that $W(t,x)=w(t,x)$. Let $\epsilon>0$ and deifne $w^{\epsilon}=w+\epsilon\xi$ where
        
        \[\xi(s,y)=e^{-(\abs{s-t}^2+\abs{y-x}^2)}-1\]

        defined over $\overline{Q}$. Since $\mathcal{D}$ is a vector space and $\epsilon\xi\in\mathcal{D}$ by \ref{4-1-condonCandD}, then 
        $w^{\epsilon}\in\mathcal{D}$. Then for every $\epsilon>0$ the function $W-w^{\epsilon}$ has a strict minimum at $(t,x)$ with $W(t,x)=w^{\epsilon}(t,x)$ which 
        implies that
        
        \[-\frac{\partial}{\partial t}w^{\epsilon}(t,x) + F(t,x,D_xw^{\epsilon}(t,x),D_x^2w^{\epsilon}(t,x),w^{\epsilon}(t,x))\leq0.\]

        By continuity of $F$ and smoothness of $w^{\epsilon}$ we let $\epsilon\to0$ and get the wanted inequality.
    \end{proof}
\end{lemma}


Finally, we prove the equivalence of these definitions. 

\begin{theorem}
    Let all the previous assumptions and $W\in C_p(\overline{Q})\cap\mathcal{M}(\overline{Q})$ and $\mathcal{D}\subset C^{1,2}(Q)$. Then $W$ is a subsolution (supersolution) 
    of \ref{4-1-PDEwithW} as defined by \ref{4-1-defviscsol} if and only if it is a subsolution (supersolution) as defined by \ref{4-1-defviscforW}. 

    \begin{proof}
        Let $W$ be a supersolution as in \ref{4-1-defviscsol}. We just prove the supersolution property, since if we have $W$ subsolution we can see it as 
        a supersolution. Let $W$ a subsolution and $(t,x)$ a strict maximum of $W-w$ and $W(t,x)=w(t,x)$, then we can find an open set $A$ such that $w>W$ over 
        $A\cap Q$. By sub-polynomial growth there exists a constant $K$ such that 

        \[\abs{W(s,y)} \leq K(1+\abs{y}^{2m}),\,\forall (s,y)\in\overline{Q}\]

        for some $m>0$. Then for some $\xi$ given by Urysohn’s lemma we can define 

        \[\overline{w}(s,y)= \xi(s,y)w(s,y) - (1-\xi(s,y))K(1+\abs{y}^{2m}),\]

        and get a minimum for $W-\overline{w}$ in $(t,x)$. Then we are in the supersolution case.

        Let $w\in C^{\infty}(Q)$ and $(t,x)$ a local minimum of $W-w$ where it is null. There exists an open subset $\mathcal{N}$ of $\R^n$ such that 
        $W\geq w$ over $\mathcal{N}\cap\overline{Q}$, and by Urysohn’s lemma there exists a function $\xi\in C^{\infty}(\overline{Q})$ taking values in $[0,1]$, 
        constantly $1$ in a neighborhood $\tilde{\mathcal{N}}$ of $(t,x)$ and zero outside $\mathcal{N}$. Since $W$ is bounded below we define 
        
        \begin{equation}
            \tilde{w}(s,y) = \xi(s,y)w(s,y) - (1-\xi(s,y))K,
        \end{equation}

        where $-K$ is the lower bound of $W$. Then, by hypotheses on $W$, $w$, and $\xi$ we have

        \[W = \xi W  + (1-\xi) W \geq \xi w + (1-\xi)W\geq\xi w-(1-\xi)K=\overline{w}.\]

        Furthermore, $\overline{w}\in\mathcal{D}$ since \ref{4-1-condonCandD} holds, $(t,x)$ is a minimum of $W-\overline{w}$, and $W(t,x)=\overline{w}(t,x)$ 
        then 
        
        \[-\frac{\partial}{\partial t}\overline{w}(t,x) + F(t,x,D_y\overline{w}(t,x),D_y^2\overline{w}(t,x),\overline{w}(t,x))\geq0,\]

        but $\overline{w}$ and $w$ have same derivative in $(t,x)$, since $\xi$ is constantly unitary in a neighborhood of it. Then we get

        \[-\frac{\partial}{\partial t}w(t,x) + F(t,x,D_yw(t,x),D_y^2w(t,x),w(t,x))\geq0.\]

        Let us suppose that $W$ is a subsolution for \ref{4-1-defviscforW}. Let $w\in\mathcal{D}$ and $(t,x)$ a strict maximum of $W-w$ such that $(W-w)(t,x)=0$. 
        Since $\mathcal{D}\subset C^{1,2}(Q)$ we can find an open set $Q^{\ast}$ over which $w$ is $C^{1,2}$ and then extend it to be null over $\R^{n+1}\setminus Q^{\ast}$. 
        Let us define a sequence of mollifications of $w$ defined as $w^n$ such that $w^n,w_t^n,w_{x_i}^n,w_{x_ix_j}^n$ converge to their respective of $w$. 
        Since the convergence of $w^n$ to $w$ is uniform there must be a subsequence $(t_n,x_n)$ converging to $(t,x)$, such that $W-w^n$ has a maximum. If we set 
        $\overline{w}^n=w^n-w^n(t_n,x_n)+W(t_n,x_n)$ we still have a maximum in $(t_n,x_n)$ for $W-\overline{w}^n$ and $W(t_n,x_n)=\overline{w}^n(t_n,x_n)$. Then 

        \[-\frac{\partial}{\partial t}\overline{w}^n(t_n,x_n) + F(t_n,x_n,D_x\overline{w}^n(t_n,x_n), D_x^2\overline{w}^n(t_n,x_n), \overline{w}^n(t_n,x_n))\leq0.\]

        Because of continuity of $F$ and smoothness of $\overline{w}^n$ we get
        
        \[-\frac{\partial}{\partial t}w(t,x) + F(t,x,D_yw(t,x),D_y^2w(t,x),w(t,x))\leq0.\]
    \end{proof}
\end{theorem}

\subsection{Value Function as Viscosity Solution}

In the first section we proved that the value function of a dynamic programming problem is the classical solution of the dynamic programming equation, under differentiability 
assumptions. We now prove that it is a viscosity solution of the dynamic programming equation given two sets of assumptions.

\begin{theorem}\label{4-1-viscodynamiceq}
    Let $U$ be a bounded space of control and $f\in C(\overline{Q}\times U)$ such that $\abs{f(t,x,v)}\leq K(1+\abs{x})$. The for every $w\in C^1(Q)\cap \mathcal{M}(\overline{Q})$ 

    \begin{equation}
        \lim_{h\to 0} \frac{1}{h}\left[(\mathcal{T}_{tt+h}w(t+h,\cdot))(x)-w(t,x)\right] = \frac{\partial}{\partial t}w(t,x) - H(t,x,D_xw(t,x)), \,\forall(t,x)\in\overline{Q}.
    \end{equation}
\end{theorem}

\vspace{5mm}


Note that equation \ref{4-1-viscodynamiceq} implies that $V$ is a viscosity solution of \ref{1-2-HJB1} via theorem \ref{4-1-thm:Vtoviscosity}.

\vspace{5mm}

\begin{proof}

    Let $(t,x)\in Q$, $v\in U$ and $w\in C^1(Q)\cap\mathcal{M}(\overline{Q})$. We prove that
    
    \begin{equation}\label{4-1-inproof:limsup}\limsup_{h\to0}\frac{1}{h}\left[(\mathcal{T}_{tt+h}w(t+h,\cdot))(x)-w(t,x)\right]\leq\frac{\partial}{\partial t}w(t,x) - H(t,x,D_xw(t,x)),\end{equation}

    and 

    \begin{equation}\label{4-1-inproof:liminf}\liminf_{h\to0}\frac{1}{h}\left[(\mathcal{T}_{tt+h}w(t+h,\cdot))(x)-w(t,x)\right]\geq\frac{\partial}{\partial t}w(t,x) - H(t,x,D_xw(t,x)).\end{equation}
    
    Equation \ref{4-1-inproof:limsup} can be proven without assuming boundedness of the control space $U$ and sub-linearity of the flow $f$. Indeed, there exists an admissible control $u(\cdot)\in \mathcal{U}(t,x)$ 
    such that $\lim_{s\to t}u(s)=v$. Let $x(\cdot)$ be the state solution of the Cauchy problem associated 
    to the control $u(\cdot)$ such that $x(t)=x$. Let also $\tau$ the exit time of $(s,x(s))$ from $Q$. If we take a sufficiently small $h>0$ we can compute
    
    \begin{align*}
        \frac{1}{h}\left[(\mathcal{T}_{tt+h}w(t+h,\cdot))(x)-w(t,x)\right] & \leq \frac{1}{h}\int_t^{t+h}L(s,x(s),u(s))\,ds + \frac{1}{h}\left[w(t+h,x(t+h))-w(t,x)\right] \\
        & =  \frac{1}{h}\int_t^{t+h}\bigg[L(s,x(s),u(s)) \\ 
        & \quad + \frac{\partial}{\partial t}w(s,x(s)) + f(s,x(s),u(s))\cdot D_xw(s,x(s))\bigg]\,ds.
    \end{align*}
    
    Then \ref{4-1-inproof:limsup} holds true. 
    
    Let us prove \ref{4-1-inproof:liminf}. Using $\abs{f(t,x,v)}\leq K(1+\abs{x})$ we get

    \begin{align*}
        \abs{x(r)-x} & \leq \int_r^t \abs{f(s,x(s),u(s))} \,ds \\
        & \leq K(r-t) + K\int_t^r \abs{x}(s)\,ds \\
        & \leq K(r-t)(1+\abs{x}) + K\int_t^r \abs{x(s)-x}\,ds,
    \end{align*}

    which implies by Gronwall's inequality 

    \[\abs{x(r)-x} \leq (1+\abs{x})e^{K(t-r)-1},\,\forall r\geq t.\]

    Then for suitably small $r$ that $\tau\geq r\wedge t_1$ for all $u(\cdot)\in\mathcal{U}(t,x)$. Then for $n$ sufficiently large by definition of $\mathcal{T}_{tr}$ we can find a control $u^n(\cdot)$ such that
    
    \begin{align*}
        \left(\mathcal{T}_{tt+\frac{1}{n}}w\left(t+\frac{1}{n}\right)\right)(x) \geq \int_t^{t+1/n} L(s,x(s),u(s)) \,ds + w\left(t+\frac{1}{n},x^n\left(t+\frac{1}{n}\right)\right) - \frac{1}{n^2},
    \end{align*}

    which implies that 

    \begin{equation}\label{4-1-inproof:geq}
        \begin{aligned}
            n\left[\left(\mathcal{T}_{tt+\frac{1}{n}}w\left(t+\frac{1}{n}\right)\right)(x) - w(t,x)\right] & \geq \frac{\partial}{\partial t}w(t,x) + n\int_t^{t+1/n} L(s,x(s),u^n(s))\,ds \\ 
            & + n\int_t^{t+1/n} f(t,x,u^n(s))\,ds \cdot D_xw(t,x) + e(n),
        \end{aligned}        
    \end{equation}
    
    where 

    \begin{equation}
        \begin{aligned}
            e(n) = & -\frac{1}{n} +n\int_t^{t+1/n} \frac{\partial}{\partial t}\big(w(s,x^n(s)) - w(t,x)\big)\,ds \\
            & + n\int_t^{t+1/n} L((s,x^n(s),u^n(s))) - L(t,x,u^n(s))\,ds \\
            & + n\int_t^{t+1/n} f(s,x^n(s),u^n(s))\cdot D_xw(s,x^n(s)) - f(t,x,u^n(s))\cdot D_xw(t,x) \,ds.
        \end{aligned}
    \end{equation}

    Because of the continuity of $L$, of $f$ and smoothness of $w$ we have that $\lim_{n\to+\infty}e(n)=0$. If we prove that

    \begin{equation}\label{miao}
        n\int_t^{t+1/n} L(s,x(s),u^n(s))\,ds + n\int_t^{t+1/n} f(t,x,u^n(s))\,ds \cdot D_x w(t,x) \leq \sup_{v\in U}\left\{f(t,x,v)\cdot p + L(t,x,v)\right\},
    \end{equation}  

    then from \ref{4-1-inproof:geq} we get

    \[\liminf_{h\to0}\frac{1}{h}\left[\left(\mathcal{T}_{tt+h}w\left(t+h\right)\right)(x) - w(t,x)\right] \geq \frac{\partial}{\partial t}w(t,x) - H(t,x,D_x w(t,x)).\]
 
    \color{red}{PROVE THAT \ref{miao} VIA CONVEXHULL STUFF.}\color{black}{}
\end{proof}

It is important to notice that the previous result holds under quite stringent hypotheses. We both require the space control to be bounded and the flow $f$ to have sublinear growth in its state variable. We can relax these assumptions 
and get a similar result, but we need to require the existence of an optimal control.\

\begin{theorem}
    If for each $(t,x)\in Q$ there exists a $u^{\ast}\mathcal{U}(t,x)$ be an optimal control, then a continuous value function is a viscosity solution of its dynamic programming equation.
    
    \begin{proof}
        Let us drop the $^{\ast}$. As we did before in theorem \ref{4-1-viscodynamiceq}, for every $w\in C^{\infty}(Q)$ we have

        \[\limsup_{h\to0}\frac{1}{h}\left[(\mathcal{T}_{tt+h}w(t+h,\cdot))(x)-w(t,x)\right]\leq\frac{\partial}{\partial t}w(t,x) - H(t,x,D_xw(t,x)).\]

        If $V-w$ has a maximum at $(t,x)$ and $V(t,x)=w(t,x)$ then 

        \[\limsup_{h\to0}\frac{1}{h}\left[\underbrace{\int_t^{t_1} L(s,x(s),u(s))\,ds}_{=V(t,x)-w(t,x)=0} + w(t+h,x) - V(t+h,x)\right]=0,\]

        which proves subsolution. Let the subsolution hypotheses hold, then from the dynamic programming principle we get

        \begin{align*}
            0 & \geq \int_t^{t+h} L(s,x(s),u(s)) \,ds + w(t+h,x(t+h)) - w(t,x)\\ 
            & \geq  \int_t^{t+h} L(s,x(s),u(s)) + \frac{\partial}{\partial t}w(s,x(s)) + f(s,x(s),u(s))\cdot D_x w(s,x(s)) \,ds \\
            & \geq \int_t^{t+h} L(s,x(s),u(s)) - H(s,x(s),D_x w(s,x(s))) \,ds.
        \end{align*} 

        Finally, if we let $h\to0$ we get

        \[-\frac{\partial }{\partial t} w(t,x) + H(t,x,D_xw(t,x)) \geq 0.\]

    \end{proof}
\end{theorem}

We now turn our attention to uniqueness property of viscosity solution to the dynamic programming equation. Let us consider the dynamic programming equation 
in a more general form

\begin{equation}\label{4-2-firstordereq}
    -\frac{\partial}{\partial t}V(t,x) + H(t,x,D_xV(t,x)) = 0,\,(t,x)\in Q.
\end{equation}

We will show that under technical assumptions on $H$, which are satisfied by the Hamiltonian, the supremum norm of the difference between two viscosity solution 
can be computed only on the frontier. Therefore, if we impose boundary conditions we get uniqueness of the viscosity solution. 
The assumptions on $H$ are

\begin{itemize}\label{4-2-Hassumptions}
    \item There exists $K\in\R$ and $h\in C([0,\infty))$ with $h(0)=0$ such that $\forall (t,x),(s,y)\in\overline{Q}$ and $\forall p,p'\in\R^n$ we have
    
    \[H(t,x,p)-H(s,y,p')\leq h(\abs{t-s}+\abs{x-y})+h(\abs{t-s})\abs{p}+K\abs{x-y}\abs{p} + K\abs{p-p'}.\]
    
    \item There exists $K'\in\R$ such that
    
    \[\abs{H_t}+\abs{H_x}\leq K'(1+\abs{p}),\]

    and

    \[\abs{H_p}\leq K'.\]
\end{itemize}

Under this hypotheses we can state the following theorem.

\begin{theorem}\label{4-2-Theo:Suponbound}
    Let $W$ and $V$ viscosity subsolution and supersolution of \ref{4-2-firstordereq} in $Q$, respectively. If $Q$ is unbounded we assume $W$, $V$ to be bounded and 
    uniformly continuous on its closure. Then
    
    \[\sup_{\overline{Q}}[W-V] = \sup_{\partial^{\ast}Q}[W-V].\]
\end{theorem}

This Theorem implies that if $V$ and $W$ are two viscosity solution of \ref{4-2-firstordereq} satisfying the boundary conditions

\begin{equation}\label{4-2-TerminalcondV}
    \begin{aligned}
        V(t,x) = g(t,x),&\forall(t,x)\in[T_0,t_1)\times O, \\
        V(t_1,x) = \psi(x),&x\in\overline{O},
    \end{aligned}
\end{equation}

then, if the assumption of theorem \ref{4-2-Theo:Suponbound} are satisfied, we get

\[\norm{W-V}_{\infty} = \sup_{\overline{Q}}[W-V]=\sup_{\partial^{\ast}Q}[W-V] = 0.\]

The proof of Theorem \ref{4-2-Theo:Suponbound} is rather long; we first prove the $\overline{Q}$ bounded case and then the unbounded one.

\begin{proof}
    
    Let $Q$ be bounded. For $\epsilon, \delta, \beta >0$ define
    
    \begin{equation}\label{4-2-InProofSup:phi}
        \Phi(t,x,s,y) = W(t,x) - V(s,y) - \frac{1}{2\epsilon} \abs{x-y}^2 -\frac{1}{2\delta} \abs{t-s}^2 + \beta(s-t_1).   
    \end{equation}

    Since $\overline{Q}$ is bounded the function $\Phi$ achieves its supremum at $(t',x'),(s',y')\in\overline{Q}\times\overline{Q}$. We will essentially prove 
    that $\Phi$ divides $W-V$ from its supremum over the boundary in an inequality chain.
    
    \begin{enumerate}
        \item Let $\rho>$ and define
        
        \begin{equation}\label{4-2-InProofSup:Drho}
            D_{\rho} = \left\{(t,x),(s,y)\in\overline{Q}\times\overline{Q}\,|\,\abs{t-s}^2+\abs{x-y}^2\leq\rho\right\},
        \end{equation}

        and 

        \[m_W(\rho)=2\left\{\abs{W(t,x)-W(s,y)}:\,((t,x),(s,y))\in D_{\rho}\right\},\]

        \[m_V(\rho)=2\left\{\abs{V(t,x)-V(s,y)}:\,((t,x),(s,y))\in D_{\rho}\right\},\]

        and 

        \[K_1 = \sup\{m_w(\rho):\,\rho\geq0\}.\]

        Since $\overline{Q}$ is compact and $W$,$V$ are continuous they are uniformly continuous, therefore $m_W,m_V\in C([0,\infty))$. Clearly 
        $m_W(0)=m_V(0)=0$. We now claim two estimates on $\abs{x-y}$ and $\abs{t-s}$:

        \begin{equation}
            \begin{aligned}
                \abs{t'-s'}\leq\sqrt{K_1\delta} \\
                \abs{x'-y'}\leq \sqrt{\epsilon m_W(K_1[\epsilon+\delta])}.
            \end{aligned}
        \end{equation}

        From the maximum in $((t',x'),(s',y'))$ we have

        \[\Phi(t',x',s',y') \leq \Phi(s',y',s',y'),\]

        implies

        \[\frac{1}{\epsilon}\abs{x'-y'}^2 + \frac{1}{\delta}\abs{t'-s'}^2\leq 2(W(t',x')-W(s',y')) \leq m_W(\abs{t'-s'}+\abs{x'-y'}).\]

        Then we have

        \[\frac{1}{\delta}\abs{t'-s'}^2\leq K_1,\,\frac{1}{\epsilon}\abs{x'-y'}^2\leq K_1,\]

        which also imply

        \[\frac{1}{\epsilon}\abs{x'-y'}^2\leq m_w(K_1[\epsilon+\delta]).\]

        \item We now show that if $(t',x')\in\partial^{\ast}Q$ or $(s',y')\in\partial^{\ast}Q$ we have
        
        \begin{equation}\label{4-2-InProofSup:phiestim1}
            \Phi(t',x',s',y') \leq\frac{1}{2}m_V(K_1(\epsilon+\delta)) + \sup_{\partial^{\ast}Q}[W-V],
        \end{equation}

        or 

        \begin{equation}\label{4-2-InProofSup:phiestim2}
            \Phi(t',x',s',y') \leq\frac{1}{2}m_W(K_1(\epsilon+\delta)) + \sup_{\partial^{\ast}Q}[W-V],
        \end{equation}

        respectively. Indeed, in the first case, we have

        \begin{align*}\Phi(t',x',s',y')&\leq W(t',x') - W(s',y') \\ & \leq  V(t',x') - V(s',y') + \sup_{\partial^{\ast}Q}[W-V] \\
        & \leq \frac{1}{2}m_W(\abs{t'-s'}^2+\abs{x'-y'}^2) +  \sup_{\partial^{\ast}Q}[W-V]. \end{align*}

        Same reasoning for the latter.

        \item We now construct an estimate over the parameter $\beta$ if $(t',x'),(s',y')\in Q$. Let us consider the test function
        
        \[w(t,x) = \frac{1}{2\delta}\abs{t-s'}^2 + \frac{1}{2\epsilon}\abs{x-y'}^2,\]

        then $(t',x')$ is where $W-w$ is maximized and $W(t',x')=w(t',x')$, then

        \begin{equation}\label{4-2-InProofSup:qandp1}-\underbrace{\frac{1}{\delta}(t'-s')}_{\overset{def}{=}q_{\delta}} + H(t',x',\underbrace{\frac{1}{\epsilon}(x'-y')}_{\overset{def}{=}p_{\epsilon}})\leq0.\end{equation}

        Analogously, defining 

        \[w(s,y) = -\frac{1}{2\delta}\abs{t'-s}^2 - \frac{1}{2\epsilon}\abs{x'-y}^2 + \beta(s-t_1),\]

        we get:

        \begin{equation}\label{4-2-InProofSup:qandp2}
            -\beta-q_{\delta}+H(s',y',p_{\epsilon})\geq0.
        \end{equation}

        Combining \ref{4-2-InProofSup:qandp1} and \ref{4-2-InProofSup:qandp2} we get

        \begin{align*}
            \beta & \leq H(s',y',p_{\epsilon}) - H(t',x',p_{\epsilon}) \\
            & \leq h(\abs{t'-s'} + \abs{x'-y'}) + h(\abs{t'-s'})\abs{p_{\epsilon}} + K\abs{x'-y'}\abs{p_{\epsilon}} \\
            & \leq h(\sqrt{K_1\epsilon} + \sqrt{K_1\delta}) + h(\sqrt{K_1\delta})\frac{1}{\epsilon}\sqrt{K_1\epsilon} + K\frac{\abs{x'-y'}^2}{\epsilon} \\
            & \leq h(\sqrt{K_1\epsilon} + \sqrt{K_1\delta}) + h(\sqrt{K_1\delta})\sqrt{\frac{K_1\epsilon}{\epsilon}} + Km_W(K_1[\epsilon+\delta]).
        \end{align*}
        
        \item We denote
        
        \begin{equation}\label{4-2-inproofSup:k}
            k(\epsilon,\delta) = h(\sqrt{K_1\epsilon} + \sqrt{K_1\delta}) + h(\sqrt{K_1\delta})\sqrt{\frac{K_1\epsilon}{\epsilon}} + Km_W(K_1[\epsilon+\delta]),
        \end{equation}

        then

        \[\lim_{\epsilon\to0}\lim_{\delta\to0} k(\epsilon,\delta) = 0.\]

        Then for suitably small $\epsilon$ and $\delta$ we have

        \[k(\epsilon,\delta) < \beta,\]

        which implies that either $(t',x')$ or $(s',y')$ is in $\partial^{\ast}Q$. Then by \ref{4-2-InProofSup:phiestim1} 
        or \ref{4-2-InProofSup:phiestim2} we have

        \begin{equation}\label{4-2-InProofSup:limsuplimsup}
            \limsup_{\epsilon\to0}\limsup_{\delta\to0} \Phi(t',x',s',y') \leq \sup_{\partial^{\ast}Q}[W-V],
        \end{equation}

        for every $\beta>0$. Finally, we have

        \[W(t,x) - V(s,y) + \beta(t-t_1) = \Phi(t,x,t,x) \leq \Phi(t',x',s',y'),\]

        then by \ref{4-2-InProofSup:limsuplimsup} we get

        \[W(t,x) - V(s,y) = \lim_{\beta\to0}W(t,x) - V(s,y) + \beta(t-t_1)\leq \lim_{\beta\to0}\Phi(t',x',s',y')\leq \sup_{\partial^{\ast}Q}[W-V].\]
    \end{enumerate}

    Let us suppose that $Q$ is unbounded, and $V,W$ are bounded and uniformly continuous. We keep the same notation as before. For every $\gamma>0$ 
    there exist $(t_{\gamma},x_{\gamma}),(s_{\gamma},y_{\gamma})\in \overline{Q}$ such that
    
    \begin{equation}
        \Phi(t_{\gamma},x_{\gamma},s_{\gamma},y_{\gamma}) \geq \sup_{\overline{Q}\times\overline{Q}}\Phi - \gamma,
    \end{equation}

    since $W$ and $V$ are bounded. We therefore define

    \begin{equation}
        \Phi_{\gamma}(t,x,s,y) = \Phi(t,x,s,y) - \frac{\gamma}{2}\left[\abs{t-t_{\gamma}}^2+\abs{s-s_{\gamma}}^2 + \abs{x-x_{\gamma}}^2 + \abs{y-y_{\gamma}}^2\right],
    \end{equation}

    for $(t,x),(s,y)\in \overline{Q}$. 
    
    We now follow a similar procedure to the one used in the bounded case.

    \begin{enumerate}[label=\arabic*')]
        \item Let the moduli of continuity $m_W,m_V$ and the bound $K_1$ as before. Since we assume the value functions to be bounded and 
        absolutely continuous we have, as before, that
        
        \[m_W,m_V\in C([0,\infty)),\, m_W(0)=m_V(0)=0,\, K_1=\sup\{m_W(\rho):\,\rho\geq0\}<+\infty.\]

        Let us note that for $(t,x),(s,y)\in\overline{Q}\times\overline{Q}$ such that 

        \[\abs{t-t_{\gamma}}^2+\abs{s-s_{\gamma}}^2 + \abs{x-x_{\gamma}}^2 + \abs{y-y_{\gamma}}^2>2,\]

        then

        \[\Phi_{\gamma}(t,x,s,y) \leq \Phi(t,x,s,y) - \gamma \leq \Phi(t_{\gamma},x_{\gamma},s_{\gamma},y_{\gamma}) = \Phi_{\gamma}\Phi(t_{\gamma},x_{\gamma},s_{\gamma},y_{\gamma}).\]

        The supremum of $\Phi_{\gamma}$ is achieved since $\Phi_{\gamma}$ is bounded by $\Phi(t_{\gamma},x_{\gamma},s_{\gamma},y_{\gamma})$. It is achieved at $(t',x'),(s',y')$. Then

        \begin{equation}
            \abs{t-t_{\gamma}}^2+\abs{s-s_{\gamma}}^2 + \abs{x-x_{\gamma}}^2 + \abs{y-y_{\gamma}}^2\leq2.
        \end{equation}

        We now claim two estimates on the distance between the time variables and space variables. For $\gamma,\delta,\epsilon\leq\frac{1}{2}$ we have

        \begin{equation}\label{4-2-inproofSup:unboundedCLAIMS}
            \abs{t'-s'} \leq \sqrt{2(K_1+1)\delta},\,\abs{x'-y'}\leq\sqrt{2\epsilon(2\gamma + m_W(2\epsilon(K_1+1)))}.
        \end{equation}

        Indeed, $\Phi_{\gamma}(t',x',s',y')\geq\Phi_{\gamma}(s',x',s',y')$ implies

        \begin{align*}
            \frac{1}{2\delta}\abs{t'-s'}^2 & \leq W(t',x')-W(s',x') + \frac{\gamma}{2}[\abs{s'-t_{\gamma}}^2-\abs{t'-t_{\gamma}}^2] \\
            & \leq \frac{1}{2}K_1 + \gamma\abs{t'-s'}^2 + \frac{\gamma}{2}\abs{t'-t_{\gamma}}^2 \\
            & \leq \frac{1}{2}K_1 + \gamma\abs{t'-s'}^2 + \gamma,
        \end{align*}

        which implies, recalling the bound on $\delta$ and $\gamma$, that 

        \[\left(\frac{1}{4\delta}\right)\abs{t'-s'}^2\leq\left(\frac{1-\delta}{2\delta}\right)\abs{t'-s'}^2\leq\left(\frac{1}{2\delta}-\gamma\right)\abs{t'-s'}^2\leq \frac{1}{2}K_1 + \frac{1}{2}.\]

        From $\Phi_{\gamma}(t',x',s',y')\geq\Phi_{\gamma}(t',y',s',y')$ we get

        \begin{align*}
            \frac{1}{2\epsilon}\abs{x'-y'}^2 & \leq W(t',x')-W(t',x') + \frac{\gamma}{2}[\abs{y'-x_{\gamma}}^2-\abs{x'-x_{\gamma}}^2] \\
            & \leq \frac{1}{2}m_W(\abs{x'-y'}^2) + \gamma\abs{x'-y'}^2 + \frac{\gamma}{2}\abs{x'-x_{\gamma}}^2 \\
            & \leq \frac{1}{2}m_W(\abs{x'-y'}^2) + \gamma\abs{x'-y'}^2 + \gamma,
        \end{align*}

        which implies

        \[\frac{1}{4\epsilon}\abs{x'-y'}^2\leq \frac{1}{2}m_W(\abs{x'-y'}^2)+\gamma\Rightarrow \abs{x'-y'}^2\leq 2(K_1+1)\epsilon.\]

        Therefore, we get

        \[\frac{1}{4\epsilon}\abs{x'-y'}^2\leq \left(\frac{1}{2\epsilon}-\gamma\right)\abs{x'-y'}^2\leq \frac{1}{2}m_W\left(2(K_1+1)\epsilon\right)+\gamma,\]

        hence the second equation in \ref{4-2-inproofSup:unboundedCLAIMS}.

        \item Replicating the same argument of the bounded case we get if $(t',x')\in\partial^{\ast}Q$ then
        
        \begin{equation}
            \Phi_{\gamma}(t',x',s',y') \leq \frac{1}{2}m_V\left(2(K_1+1)(\epsilon+\delta)\right) + \sup_{\partial^{\ast}Q}[W-V],
        \end{equation}

        and if $(s',y')\in\partial^{\ast}Q$ then

        \begin{equation}
            \Phi_{\gamma}(t',x',s',y') \leq \frac{1}{2}m_W\left(2(K_1+1)(\epsilon+\delta)\right) + \sup_{\partial^{\ast}Q}[W-V].
        \end{equation}

        \item Finally, if we find a bound over $\beta$ for $(t',x'),(s',y')\in Q$ we then prove that this bound is not satisfied and then we can invoke 
        point $2'$ and get

        \begin{equation}
            \limsup_{\epsilon\to0}\limsup_{\delta\to0}\limsup_{\gamma\to0}\Phi_{\gamma}(t',x',s',y') \leq \sup_{\partial^{\ast}Q}[W-V],
        \end{equation}

        which implies, as before, that

        \[W(t,x) - V(s,y)\leq \sup_{\partial^{\ast}Q}[W-V].\]

        Let $(t',x'),(s',y')\in Q$ and the test function

        \[\overline{w}(t,x)=\frac{1}{2\delta}\abs{t-s'}^2+\frac{1}{2\epsilon}\abs{x-y'}^2+\frac{\gamma}{2}\left(\abs{t-t_{\gamma}}^2+\abs{x-x_{\gamma}}^2\right),\]

        defined over $Q$. Then $(t',x')$ is where $W-\overline{w}$ is null and maximized. Therefore

        \begin{equation}\label{4-2-inproofSup:qandpwithoutbeta}
             - \underbrace{\frac{1}{\delta}(t'-s')}_{\overset{def}{=}q_{\delta}}  - \underbrace{\gamma(t'-t_{\gamma})}_{\overset{def}{=}q_{\gamma}} + H(t',x',p_{\epsilon}+p_{\gamma}) \leq 0, 
        \end{equation}

        where

        \[p_{\epsilon}=\frac{1}{\epsilon}(x'-y'),\, p_{\gamma}=\gamma(x'-x_{\gamma}).\]

        With the same reasoning we define the test function

        \[w^{\ast}(s,y)=-\frac{1}{2\delta}\abs{t'-s}^2-\frac{1}{2\epsilon}\abs{x'-y}^2+\frac{\gamma}{2}\left(\abs{s-s_{\gamma}}^2+\abs{y-y_{\gamma}}^2\right)+\beta(s-t_1),\]

        from which we get

        \begin{equation}\label{4-2-inproofSup:qandpwithbeta}
            - \beta - q_{\delta}  - \underbrace{\gamma(s'-s_{\gamma})}_{\overset{def}{=}\overline{q}_{\gamma}} + H(s',y',p_{\epsilon}+\overline{p}_{\gamma}) \geq 0, 
       \end{equation}

       where $\overline{p}_{\gamma} = \gamma(y_{\gamma}-y').$

       We can find the upper bound to $\beta$ using the estimates in \ref{4-2-inproofSup:unboundedCLAIMS} and equations \ref{4-2-inproofSup:qandpwithbeta} and \ref{4-2-inproofSup:qandpwithoutbeta}:

       \begin{align*}
       \beta&\leq H(s',y',p_{\epsilon}+\overline{p}_{\gamma}) - H(t',x',p_{\epsilon}+p_{\gamma}) + q_{\gamma} - \overline{q}_{\gamma} \\
       & \leq h(\abs{t'-s'}+\abs{x'-y'})+h(\abs{t'-s'})\left[\abs{p_{\epsilon}+\abs{\overline{p}_{\gamma}}}\right] \\
       & + K\abs{x'-y'}\left[\abs{p_{\epsilon}+\abs{\overline{p}_{\gamma}}}\right]+K\abs{\overline{p}_{\gamma}-p_{\gamma}}+q_{\gamma}-\overline{q}_{\gamma} \\
       & \leq h(\sqrt{2(K_1+1)\delta}+\sqrt{2(K_1+1)\epsilon})+h(\sqrt{2(K_1+1)\delta})\left[\sqrt{\frac{\abs{2(K_1+1)}}{\epsilon}}+2\gamma\right] \\
       & + K2\left[2\gamma+m_W(2(K_1+1))\epsilon+\sqrt{2(K_1+1)\epsilon\delta}\right]+2(K+1)\gamma.
       \end{align*}
    \end{enumerate}
\end{proof}

\subsection{Continuity Properties of the Value Function}

We now investigate continuity properties of the value function. In particular, we consider two different sets of assumptions which guarantee continuity, actually 
something stronger: Lipschitz continuity and Uniform continuity. Let us consider $f,L,\psi$ as usual, recalling that there exists $K_{\rho}>0$ such that 

\begin{equation}\label{4-3-boundonf}
    \abs{f(t,x,v) - f(t,y,v)} \leq K_{\rho}\abs{x-y}, 
\end{equation}

for all $\abs{v}\leq\rho$. Note that these are minimal assumption to start working with the optimal control problem.

\begin{theorem}\label{4-3-Theorem:1cont}
    Let a bounded control space $U$, $Q=[t_0,t_1)\times\R^n$. Assume that $f,L,\psi$ are bounded, $f$ satisfies \ref{4-3-boundonf} and $L,\psi$ uniformly continuous. 
    Then the value function $V$ is bounded and uniformly continuous.
    
    \begin{proof}
        $V$ is bounded. Indeed 

        \[V(t,x) = \inf_{u\in\mathcal{U}}\int_t^{t_1} L(s,x(s),u(s))\,ds + \psi(t_1,x(t_1)) \leq (t_1-t)K_L + K_{\psi},\]

        where $K_g$ is the bound over the function $g$. 

        $V$ is uniformly continuous. Standard arguments on solutions of ODE (integral form of the solution, bound on the flow and Gronwall) yield 

        \begin{equation}
            \Lambda(s) \overset{def}{=} \abs{x(s)-y(s)} \leq \abs{x-y}e^{K(s-t)},
        \end{equation} 

        where $x(s),y(s)$ are solution to the Cauchy problem

        \[\frac{d}{ds}w(s) = f(s,w(s),u(s)),\, w(t)=w.\]

        Then

        \begin{align*}
            \abs{J(t,x,u) - J(t,y,u)} & \leq \int_t^{t_1} \abs{L(s,x(s),u(s)) - L(s,y(s), u(s))}\,ds + \abs{\psi(x(t_1)) - \psi(y(t_1))} \\
            & \leq \int_t^{t_1} m_L(\Lambda(s))\,ds + m_{\psi}(\Lambda(t_1)) \\
            & + (t_1-t_0)m_L(\abs{x-y}e^{K(t_1-t_0)}) + m_{\psi}(\abs{x-y}e^{K(t_1-t_0)}),
        \end{align*}

        where $m_g$ is the modulus of continuity of the function $g$. Since both $L$ and $\psi$ are uniformly continuous $m_L,m_{\psi}\in C([0,\infty))$ and 
        they are null at the origin. Then for a fixed $t$ we have the bound
        
        \begin{equation}\label{4-3-InProof:modulust}
            \abs{V(t,x)-V(t,y)} \leq (t_1-t_0)m_L(\abs{x-y}e^{K(t_1-t_0)}) + m_{\psi}(\abs{x-y}e^{K(t_1-t_0)}).
        \end{equation}

        Similarly, if we fix $x$ we get 

        \begin{equation}\label{4-3-InProof:modulusx}
            \abs{V(t,x)-V(s,x)} \leq K_1\abs{t-s} + (t_1-t_0)m_L(\abs{t-s}K_2e^{K(t_1-t_0)}) + m_{\psi}(\abs{t-s}K_2e^{K(t_1-t_0)}),
        \end{equation}

        where $K_1=\sup\abs{L}$ and $K_2$ is such that $\abs{x(s)-x}\leq K_2\abs{t-s}$, which exists because $f$ is bounded. Then by splitting 
        $\abs{V(t,x)-V(s,y)}$ we find a modulus of continuity for it, by summing \ref{4-3-InProof:modulusx} and \ref{4-3-InProof:modulust}. 
    \end{proof}
\end{theorem}

From the proof it is clear that by taking $L$ and $\psi$ Lipschitz continuous we get $V$ Lipschitz continuous, since we can use 

\[m_L(\rho)=K_L\rho,\, m_{\psi}(\rho)=K_{\psi}\rho,\]

for some constants $K_L,K_{\psi}>0$. 

This theorem allows us to conclude the following result.

\begin{corollary} 
    Under the previous assumptions, the value function is the unique viscosity solution of the dynamic programming equation with fixed terminal conditions \ref{4-2-TerminalcondV}. 

    \begin{proof} 
        Since $U$ is bounded the Hamiltonian $H$ satisfies \ref{4-2-Hassumptions}. Indeed, the bound on the derivative is obvious. 
        The upper bound function $h$ can be found as follows:

        \begin{align*}
            H(t,x,p) - H(s,y,p') & \leq \sup_{v\in U}\left\{\abs{L(t,x,v)-L(s,y,v) + f(t,x,v)\cdot p - f(s,y,v)\cdot p'}\right\} \\
            & \leq \sup_{v\in U}\left\{\abs{L(t,x,v)-L(s,y,v)}\right\} + \sup_{v\in U}\abs{f(s,y,v)}\abs{p-p'} \\
            & + \sup_{v\in U}\abs{f(t,x,v)-f(s,y,v)}\abs{p},
        \end{align*}

        but the uniform continuity of $L$ and Lipschitz continuity of $f$ yield \ref{4-2-Hassumptions}. Then by Theorem \ref{4-2-Theo:Suponbound} we get the thesis.
    \end{proof} 
\end{corollary}

% Let us now address the bounded $Q$ case. We assume that the value function has the form

% \[V(t,x) = \inf_{u\in\mathcal{U}} J(t,x;u),\]

% where 

% \[J(t,x;u) = \int_t^{\tau} L(s,x(s),u(s)) \,ds + \Psi(\tau,x(\tau)).\]

% We require $\Psi$ to be 

% \[\Psi(\tau,x(\tau)) = 0\text{ if }\tau<t_1,\,\Psi(\tau,x(\tau)) = \psi(x(t_1))\text{ if }\tau = t_1.\]

% \begin{theorem}
    
% \end{theorem}


\subsection{Pontryagin's Principle}

\newpage

\nocite{*}
\bibstyle{numeric}
\printbibliography
\addcontentsline{toc}{chapter}{Bibliography}

% % Run this to get table of contents
% biber Main
% pdflatex Main

\end{document}